\documentclass[11pt]{amsart}

\usepackage{amsmath}
\usepackage{amssymb}
\usepackage{amsthm}
\usepackage{booktabs}
\usepackage{array}
\usepackage{hyperref}

% ---------------------------------------------------------------------------
% Theorem environments
% ---------------------------------------------------------------------------
\newtheorem{mainthm}{Theorem}
\renewcommand{\themainthm}{\Alph{mainthm}}
\newtheorem{theorem}{Theorem}[section]
\newtheorem{proposition}[theorem]{Proposition}
\newtheorem{lemma}[theorem]{Lemma}
\newtheorem{corollary}[theorem]{Corollary}
\theoremstyle{definition}
\newtheorem{definition}[theorem]{Definition}
\theoremstyle{remark}
\newtheorem{remark}[theorem]{Remark}

\numberwithin{equation}{section}
\setlength{\emergencystretch}{2em}

% ---------------------------------------------------------------------------
% Macros
% ---------------------------------------------------------------------------
\newcommand{\NN}{\mathbb{N}}
\newcommand{\ZZ}{\mathbb{Z}}
\newcommand{\QQ}{\mathbb{Q}}
\newcommand{\RR}{\mathbb{R}}
\newcommand{\CC}{\mathbb{C}}
\newcommand{\FF}{\mathbb{F}}
\newcommand{\cA}{\mathcal{A}}
\newcommand{\cR}{\mathcal{R}}
\newcommand{\eps}{\varepsilon}
\newcommand{\ind}[1]{[\,#1\,]}
\newcommand{\lev}{\operatorname{lev}}
\newcommand{\cb}{\operatorname{cb}}
\newcommand{\Li}{\operatorname{Li}}
\newcommand{\Ti}{\operatorname{Ti}}
\newcommand{\Ed}{\deg_E}
\newcommand{\odd}{\operatorname{odd}}
\newcommand{\cl}{\operatorname{cl}}
\newcommand{\ev}{\operatorname{ev}}
\newcommand{\lean}[1]{{\def\_{\textunderscore\allowbreak}\texttt{#1}}}
% Sentences that depend on Hata, J. Math. Pures Appl. 69 (1990), are marked with a comment "% HATA-n" and listed in
% HATA_TODO.md, so that they can be revised quickly if that paper is read.

\begin{document}

\hypersetup{pdftitle={Irrationality of chi2(1/q), Ti2(1/q) and chi2(2/q)},pdfauthor={Shawn Krisman}}
\title[Irrationality of $\chi_2(1/q)$, $\Ti_2(1/q)$ and $\chi_2(2/q)$]{Irrationality of $\chi_2(1/q)$, $\Ti_2(1/q)$ and
$\chi_2(2/q)$}

\author{Shawn Krisman}

\subjclass[2020]{Primary 11J72; Secondary 33B30, 41A21, 31A15, 68V20}
\keywords{Legendre chi function, inverse tangent integral, dilogarithm, irrationality, Hankel determinants, Pad\'e
approximation, logarithmic potential theory, formal verification, Lean}

\date{\today}

\begin{abstract}
Let $\chi_2(z)=\sum_{n\ge0}z^{2n+1}/(2n+1)^2$ be Legendre's chi function and
$\Ti_2(z)=\sum_{n\ge0}(-1)^nz^{2n+1}/(2n+1)^2$ the inverse tangent integral. We prove three theorems: $\chi_2(1/q)$ is
irrational for every integer $q\ge3$ with $q\ne4$; $\Ti_2(1/q)$ is irrational for every integer $q\ge5$; and
$\chi_2(2/q)$ is irrational for every odd integer $q\ge175$. All three come from one construction: the Hankel
determinants of a moment functional with half-integer poles and geometric weight $w$, read at $w=q^{-2}$, $w=-q^{-2}$
and $w=4q^{-2}$. Their denominators are bounded prime by prime, their size by a discrete log-gas argument against
explicit dual certificates, and their nonvanishing by a reduction modulo an auxiliary prime. The three cases differ at
the prime $2$, where each needs its own lemma, in the certificates, and in how the nonvanishing lemma is transferred.
% HATA-3 (begin): the novelty sentence of the abstract
For large $q$ the three statements follow from the linear independence criteria of David, Hirata-Kohno and Kawashima;
for the first two, a 1990 criterion of Hata, which we could not consult, also applies from a threshold unknown to us.
Below these ranges the results are new as far as we know.
% HATA-3 (end)
All three theorems are formalised in Lean~4 with Mathlib and checked by the Lean kernel, using only
Lean's three standard axioms.
\end{abstract}

\maketitle

\tableofcontents

% ===========================================================================
\section{Introduction}\label{sec:intro}
% ===========================================================================

\subsection{Results}
For $|z|\le1$ let
\begin{align*}
  \chi_2(z)&=\sum_{n=0}^{\infty}\frac{z^{2n+1}}{(2n+1)^2}=\tfrac12\bigl(\Li_2(z)-\Li_2(-z)\bigr),\\
  \Ti_2(z)&=\sum_{n=0}^{\infty}\frac{(-1)^nz^{2n+1}}{(2n+1)^2}=\operatorname{Im}\Li_2(iz)
\end{align*}
be Legendre's chi function and the inverse tangent integral \cite[\S1.8, Ch.~2]{Lewin1981}.

\begin{mainthm}\label{thm:A}
For every integer $q\ge3$ with $q\ne4$, the number $\chi_2(1/q)$ is irrational.
\end{mainthm}

\begin{mainthm}\label{thm:B}
For every integer $q\ge5$, the number $\Ti_2(1/q)$ is irrational.
\end{mainthm}

\begin{mainthm}\label{thm:C}
For every odd integer $q\ge175$, the number $\chi_2(2/q)$ is irrational.
\end{mainthm}

The cases $q=3$ and $q=5$ of Theorem~\ref{thm:A} were proved in \cite{Chi2Paper}, whose construction we extend here.
For even $q$, $\chi_2(2/q)=\chi_2(1/(q/2))$ is covered by Theorem~\ref{thm:A} unless $q/2\in\{1,2,4\}$, so
Theorem~\ref{thm:C} is stated for odd $q$. With finer ball-arithmetic certificates, which are not part of this paper's
proof, the argument would reach $q=173$ (Remark~\ref{rem:173}).

\subsection{Context}\label{sec:context}
In Hata's notation put $L(\alpha,k;x)=\sum_{n\ge1}x^n/(n-\alpha)^k$. Then
\begin{equation}\label{eq:lerch}
  \chi_2(\tfrac1q)=\tfrac q4L(\tfrac12,2;\tfrac1{q^2}),\quad
  \Ti_2(\tfrac1q)=-\tfrac q4L(\tfrac12,2;-\tfrac1{q^2}),\quad
  \chi_2(\tfrac2q)=\tfrac q8L(\tfrac12,2;\tfrac4{q^2}),
\end{equation}
so each of our numbers is a single shifted dilogarithm (a value of Lerch's transcendent with shift $\frac12$) at a
rational point of size about $q^{-2}$.
\begin{itemize}
\item \emph{The criterion of David, Hirata-Kohno and Kawashima.} Their Theorem~2.1 in \cite{DHK2022} (shifts $\frac12$,
  depths $\le2$, one rational point $\alpha/\beta$) gives the linear independence of $1$, $L(\frac12,1;\alpha/\beta)$,
  $L(\frac12,2;\alpha/\beta)$ over $\QQ$ when an explicit quantity $V$ is positive. By our evaluation,
  at $\alpha/\beta=\pm q^{-2}$ the condition is $\log q^2>2\log2+6\log3+8$, i.e.\ $q\ge2949$; this covers
  $\chi_2(1/q)$ and $\Ti_2(1/q)$ for $q\ge2949$. At $\alpha/\beta=4/q^2$ the criterion as printed gives $q\ge11{,}795$
  for $\chi_2(2/q)$. As far as we can see, the last step of their proof gives only the weaker form of the criterion, with
  $(\rho m+1)\log\|\alpha\|$ in place of $\rho m\log\|\alpha\|$, which gives $q\ge23{,}587$; we compare with the printed
  threshold $11{,}795$. (For odd $q$ both thresholds are the first odd integers above $11{,}793.2$ and $23{,}586.4$.) These
  specialisations are ours; \cite{DHK2022} does not mention $\chi_2$ or $\Ti_2$. The other criteria of these authors
  \cite{DHK2020,DHK2022b} give larger thresholds here.
% HATA-1 (begin)
\item \emph{Hata's criterion.} Hata \cite{Hata1990} proved linear independence measures for $1$ and values
  $L(\alpha_j,k_j;1/Q)$ with rational shifts $\alpha_j$, for integers $|Q|>Q_0$ with an explicit $Q_0$. We have not been
  able to consult this paper; what follows is secondhand, from its zbMATH review and from later citations
  \cite{DHK2020,DHK2022}. Negative $Q$ are allowed, and the stated application is $\Li_2(1/Q)\notin\QQ$ for $Q\ge12$
  or $Q\le-8$; Hata himself attributes this application to \cite{Hata1990} in \cite[\S1]{Hata1993}. Applied with $Q=q^2$ and $Q=-q^2$, the theorem therefore gives $\chi_2(1/q)\notin\QQ$ and
  $\Ti_2(1/q)\notin\QQ$ for all sufficiently large $q$, by \eqref{eq:lerch}. We do not know its threshold for the shift
  $\frac12$, at either sign of $Q$. A heuristic reconstruction of what we take to be Hata's construction at shift
  $\frac12$ and $Q=q^2$ (calibrated so that it reproduces his unshifted constant $11.837\ldots$) suggests a threshold
  near $q\approx55$--$85$ for $\chi_2(1/q)$; this is an estimate, not a statement about his paper, and it says nothing
  about $Q=-q^2$, where his unshifted thresholds are better. The point $4/q^2$ of $\chi_2(2/q)$ is not of the form
  $1/Q$, so Theorem~\ref{thm:C} lies outside the theorem as it is reviewed.
% HATA-1 (end)
\item \emph{Single dilogarithm values do not suffice.} Hata \cite[Thm.~1.1]{Hata1993} proved $\Li_2(1/k)\notin\QQ$,
  and the linear independence of $1$, $\log(1-1/k)$, $\Li_2(1/k)$, for every integer $k\ge7$ or $k\le-5$, with
  Pad\'e-type approximations to the unshifted dilogarithm built from Legendre-type polynomials; it does not treat
  shifted values. Recent work proves
  $\Li_2(1/n)\notin\QQ$ for every integer $n\ne0$ \cite{Zhang2026}, after special cases \cite{TangLi2,Prellberg2026,LiuLi2},
  and $\Li_2(p/q)\notin\QQ$ for $q>(4e^4/27)p^3$ \cite{TanZhang2026}
  (all recent, unrefereed). Two irrational numbers can have a rational difference, so none of these says anything
  about $\chi_2(z)=\frac12(\Li_2(z)-\Li_2(-z))$ or $\Ti_2(z)=\operatorname{Im}\Li_2(iz)$. The independence results of
  Rhin--Viola \cite{RhinViola2005} and Viola--Zudilin \cite{ViolaZudilin2018} concern other pairs of points.
\item \emph{Integer arguments.} Bah\c{c}{\i}van \cite{Bahcivan2026} (unrefereed) proves, from the criterion of
  \cite{DHK2022} and the reciprocity formulas, that $\Ti_s(n)$, $\operatorname{Re}\chi_2(n)$ and related values are
  irrational for a set of positive integers $n$ of natural density one. These are different numbers:
  $\Ti_2(n)=\Ti_2(1/n)+\frac\pi2\log n$ and $\operatorname{Re}\chi_2(n)=\frac{\pi^2}4-\chi_2(1/n)$ for $n>1$. His
  Corollary~6.2 ($1$, $\Ti_2(n)$, $\pi\log n$ linearly independent for a density-one set of $n$) does imply
  $\Ti_2(1/n)\notin\QQ$ on a set of density one. That set is not explicit: independence is obtained only through the
  criterion of \cite{DHK2022} at points $1/(qs)$ with $q$ sufficiently large, and no individual $n$ is decided.
\item \emph{Identities do not suffice.} Landen's identity \cite[App.~A.2.1, (24)]{Lewin1981} relates $\chi_2(x)$ to
  $\chi_2\bigl(\frac{1-x}{1+x}\bigr)$ up to a product of logarithms, and $\Ti_2(x)-\Ti_2(1/x)=\frac\pi2\log x$. Neither
  moves our numbers to values known to be irrational.
\end{itemize}
% HATA-2 (begin)
\emph{Novelty.} To our knowledge, and relative to the criteria we have read, the following ranges are new:
Theorem~\ref{thm:A} for $6\le q\le2948$ (the cases $q=3,5$ are in \cite{Chi2Paper}); Theorem~\ref{thm:B} for
$5\le q\le2948$; and Theorem~\ref{thm:C} for $175\le q<11{,}795$. For Theorems~\ref{thm:A} and~\ref{thm:B} this is
subject to the threshold of \cite{Hata1990} for the shift $\frac12$, which we could not check: if that threshold is
below $q$, the corresponding values are covered by Hata's theorem. By the heuristic estimate above, the range that is
genuinely new for Theorem~\ref{thm:A} may be much shorter than $6\le q\le2948$, perhaps $q\lesssim85$ or less; for
Theorem~\ref{thm:B} we have no estimate. Theorem~\ref{thm:C} does not depend on \cite{Hata1990} in this way.
% HATA-2 (end)
Our search for prior work covered the web, arXiv, Zenodo and GitHub (repository and code search), but not MathSciNet
or zbMATH beyond individual reviews.

\subsection{The method}\label{sec:method}
Fix a real number $w$ with $0<|w|<1$. For $h\ge1$ let $x_k:=-(k+\frac12)$ for $0\le k<h$, $D(u):=\prod_{k<h}(u-x_k)$
and
\begin{equation}\label{eq:Lambda}
  \Lambda(g):=\sum_{t\ge1}w^t\Bigl(\frac{g(u)\,u}{D(u)}\Bigr)'\Big|_{u=t}\qquad(g\in\QQ[u]).
\end{equation}
For $w=q^{-2}$, $w=-q^{-2}$ and $w=4q^{-2}$, partial fractions give $\Lambda(g)=A(g)+\theta B(g)$ with rational
functionals $A,B$ and $\theta=\chi_2(1/q)$, $\Ti_2(1/q)$ and $\chi_2(2/q)$ respectively (Lemma~\ref{lem:pole}). So the
Hankel determinant $\det(\Lambda(u^{i+j}))_{i,j<h}$ is the value at $X=\theta$ of a polynomial
$\Delta_h(X)\in\QQ[X]$ of degree at most $h$. This is the construction of \cite{Chi2Paper}, a Pad\'e-type
construction of the kind used by Q.~Tang for $\Li_2(1/2)$ \cite{TangLi2}. The proof has three inputs.
\begin{enumerate}
\item \emph{Denominators} (Sections~\ref{sec:arith1}--\ref{sec:arith2} and the prime $2$ in
  Sections~\ref{sec:A}--\ref{sec:C}). There are rationals $d_h>0$ with $d_h\Delta_h\in\ZZ[X]$ and
  $\log d_h\le(\alpha+o(1))h^2$, where $\alpha=\frac{133}{60}+c_{\mathrm{odd}}+c_2$. The term $\frac{133}{60}$ comes from
  the odd primes not dividing $q(q^2-1)$, $q(q^2+1)$ or $q(q^2-4)$ respectively; $c_{\mathrm{odd}}$ collects the odd
  primes dividing $q$ (a gain) and those dividing $q^2-1$, $q^2+1$ or $q^2-4$ (a cost); $c_2$ is the prime $2$.
\item \emph{Size} (Sections~\ref{sec:analytic}--\ref{sec:cert}). For every dual certificate,
  $\limsup h^{-2}\log|\Delta_h(\theta)|\le\Phi$, an explicit exponent that depends on $w$ only through $|w|$. A shift
  lemma moves a certificate from $|w|$ to any smaller $|w'|$ and gains $\frac14\log(|w|/|w'|)$.
\item \emph{Nonvanishing} (Section~\ref{sec:nonvan}). For every rational $a/b$, $\Delta_h(a/b)\ne0$ for infinitely
  many $h$, by a reduction modulo a prime $\ell=h/3$.
\end{enumerate}
If $\alpha+\Phi<0$, the standard argument (Proposition~\ref{prop:criterion}) concludes. Table~\ref{tab:constants}
lists the arithmetic constants.

\begin{table}[ht]
\centering\small
\begin{tabular}{@{}lllll@{}}
\toprule
 & $w$ & odd primes $c_{\mathrm{odd}}$ & prime $2$: $c_2$ & lemma\\
\midrule
\ref{thm:A}, odd $q$ & $q^{-2}$ & $\frac14\log\odd(q^2-1)-\frac13\log q$ & $\frac{\nu-2}4\log2$ & \ref{lem:twoA}\\
\ref{thm:A}, even $q$ & $q^{-2}$ & $\frac14\log(q^2-1)-\frac13\log\odd(q)$ & $(\frac12-\frac e3)\log2$ & \ref{lem:ten0}\\
 & & & or $(c(e)-\frac e3)\log2$ & \ref{lem:tene}\\
\ref{thm:B}, odd $q$ & $-q^{-2}$ & $\frac14\log\frac{q^2+1}2-\frac13\log q$ & $\frac14\log2$ & \ref{lem:ten1}\\
 & & & or $\frac18\log2$ & \ref{lem:M}\\
\ref{thm:B}, even $q$ & $-q^{-2}$ & $\frac14\log(q^2+1)-\frac13\log\odd(q)$ & $(c(e)-\frac e3)\log2$ & \ref{cor:teneTi}\\
\ref{thm:C}, odd $q$ & $4q^{-2}$ & $\frac14\log(q^2-4)-\frac13\log q$ & $\frac43\log2$ & \ref{lem:tenZ}\\
\bottomrule
\end{tabular}
\caption{The arithmetic constants per $h^2$ for Theorems~\ref{thm:A}--\ref{thm:C}, in addition to $\frac{133}{60}$.
Here $\odd(n)$ is the odd part of $n$, $\nu=v_2(q^2-1)$, $e=v_2(q)$ and $c(e)=(6-4e)/(3(3e+8))$.}\label{tab:constants}
\end{table}

The odd primes behave in the same way in all three cases, because the three determinants are specialisations of one
symbolic determinant over $\QQ[q,E]$, where $E$ stands for $1/(q^2-1)$: at $(q,E)=(q,1/(q^2-1))$ for
Theorem~\ref{thm:A}, at $(q/2,4/(q^2-4))$ for Theorem~\ref{thm:C}, and, through a coefficient twist that realises
$q\mapsto iq$, $X\mapsto-iX$, for Theorem~\ref{thm:B} (Proposition~\ref{prop:twist}). The prime $2$ is different in each
case, and needs a separate lemma in each:
\begin{itemize}
\item for odd $q$ in Theorem~\ref{thm:A}, the bound of \cite{Chi2Paper} at general $\nu=v_2(q^2-1)\ge3$;
\item for even $q$, a weighted Gauss valuation on the reflected expansion of the functional, which gives the constant
  $c(e)$; it is what brings $q=6,8,10,12$ within reach;
\item for $\Ti_2$ at odd $q$, $\nu=v_2(q^2+1)=1$, so the $2$-part of $q^2+1$ saves only $\frac14\log2$ and
  $\Ti_2(1/q)$ pays $\frac12\log2$ more in total than $\chi_2(1/q)$ (entry bounds in the anchored basis cost
  $\frac14\log2$); a change to the binomial basis halves the cost of the prime $2$ to $\frac18\log2$
  (Lemma~\ref{lem:M}), which brings $q=5$ within reach ($q=3$ stays out);
\item for $\chi_2(2/q)$ the numerator costs at least $\log2$ per $h^2$ (we prove $n_2\ge h^2-h$) and at most
  $\frac43\log2$, which the analytic side pays only from $q\approx175$ on.
\end{itemize}

\subsection{The formalisation}
All three theorems and their complete proofs are formalised in Lean~4 \cite{LeanFour} with Mathlib \cite{Mathlib} and
checked by the Lean kernel, in the repository \url{https://github.com/telaviv/chi2-irrationality} \cite{Chi2Repo},
which also contains the formalisation of \cite{Chi2Paper}. The proofs depend only on Lean's three standard axioms
(\lean{propext}, \lean{Classical.choice}, \lean{Quot.sound}); they contain no \lean{sorry} and no
\lean{native\_decide}, and every numerical certificate is evaluated by the kernel. Section~\ref{sec:lean} maps each
step to its Lean declaration. In a few places the formal proof differs from the written one (cruder lower-order terms,
and the organisation into separate files); these are listed in
Section~\ref{sec:lean-departures}.

The written proof below is self-contained, with two exceptions that are only described: the dual certificates
(Section~\ref{sec:cert}), which are decidable statements about rational data proved by evaluation in the Lean kernel,
and the closed forms of the block determinants in Section~\ref{sec:nonvan}, which are identities between explicit
rational functions.

\subsection{Organisation and notation}
Section~\ref{sec:construction} sets up the machine at a general weight and proves the pole expansion.
Section~\ref{sec:criterion} states the criterion. Sections~\ref{sec:arith1} and~\ref{sec:arith2} bound the odd part of the
denominators for all three cases at once. Sections~\ref{sec:analytic} and~\ref{sec:cert} bound the size and prove the shift
lemma, and Section~\ref{sec:nonvan} proves the nonvanishing lemma. Sections~\ref{sec:A}, \ref{sec:B} and~\ref{sec:C}
specialise to the three theorems: the prime $2$, the certificates, the nonvanishing data and the margins.
Section~\ref{sec:limits} discusses the limits of the method, and Section~\ref{sec:lean} the formalisation.

Throughout, $p$ and $\ell$ denote primes; $v_p$ is the $p$-adic valuation on $\QQ$ and $|x|_p=p^{-v_p(x)}$. For an
integer $n\ne0$, $\odd(n)=n/2^{v_2(n)}$. $\lfloor\cdot\rfloor$ and $\lceil\cdot\rceil$ are the floor and ceiling,
$(x)_+=\max(x,0)$, and $\ind{P}$ is $1$ if $P$ holds and $0$ otherwise. For a sequence $f$ on $\ZZ_{\ge0}$,
$\Delta f(t)=f(t+1)-f(t)$. $s_2(n)$ is the binary digit sum, so $v_2(n!)=n-s_2(n)$. We write $\theta(x)=\sum_{p\le
x}\log p$ only in Section~\ref{sec:arith1}; elsewhere $\theta$ is the target number.

% ===========================================================================
\section{The machine}\label{sec:construction}
% ===========================================================================

\subsection{Data}\label{sec:data}
We work with three kinds of \emph{point}. Each is given by a rational weight $w$ with $0<|w|<1$, a real target $\theta$,
and rationals $c_k$, $s_k$ ($k\ge0$):
\begin{center}
\begin{tabular}{@{}lllll@{}}
\toprule
case & $w$ & $\theta$ & $c_k$ & $s_k$\\
\midrule
(A) $q\ge2$ & $q^{-2}$ & $\chi_2(1/q)$ & $q^{2k+1}$ & $S_k(q)$\\
(B) $q\ge2$ & $-q^{-2}$ & $\Ti_2(1/q)$ & $(-1)^kq^{2k+1}$ & $S_k^{\Ti}(q)$\\
(C) $q\ge3$ odd & $4q^{-2}$ & $\chi_2(2/q)$ & $\zeta^{2k+1}$, $\zeta=q/2$ & $S_k(\zeta)$\\
\bottomrule
\end{tabular}
\end{center}
Here, for a rational $y$,
\[
  S_k(y):=\sum_{j=0}^{k}\frac{y^{2(k-j)}}{(2j+1)^2},\qquad
  S_k^{\Ti}(y):=\sum_{j=0}^{k}\frac{(-1)^{k-j}y^{2(k-j)}}{(2j+1)^2}=S_k(iy).
\]
For $h\ge1$ we use the $h$ poles $x_k=-(k+\frac12)$, $0\le k<h$, and put
\[
  D(u):=\prod_{k<h}(u-x_k),\qquad \omega_k:=\frac{x_k}{D'(x_k)},\qquad \varphi(u):=\frac u{D(u)},
\]
where $D'(x_k)=\prod_{m\ne k}(m-k)=(-1)^kk!\,(h-1-k)!$. For a rational function $F$ regular at the positive integers and of
polynomial growth put
\[
  \mu(F):=\sum_{t\ge1}w^tF'(t),
\]
an absolutely convergent series. The functional \eqref{eq:Lambda} is $\Lambda(g)=\mu(g\varphi)$.

\begin{lemma}[pole series]\label{lem:poleseries}
In each case, for every $k\ge0$,
\[
  \mu\Bigl(\frac1{u-x_k}\Bigr)=-4\bigl(c_k\,\theta-s_k\bigr).
\]
\end{lemma}

\begin{proof}
We have
\[
  \mu\Bigl(\frac1{u-x_k}\Bigr)=-\sum_{t\ge1}\frac{w^t}{(t+k+\frac12)^2}=-4\sum_{t\ge1}\frac{w^t}{(2t+2k+1)^2};
\]
put $s=2t+2k+1$, an odd number $\ge2k+3$. In case (A), $w^t=q^{-2t}=q^{2k+1}q^{-s}$, so the last sum is
\[
  q^{2k+1}\Bigl(\chi_2(1/q)-\sum_{s\le2k+1,\ s\text{ odd}}\frac{q^{-s}}{s^2}\Bigr)=q^{2k+1}\chi_2(1/q)-S_k(q).
\]
Case (C) is the same computation with $q$ replaced by $\zeta=q/2$, since $4q^{-2}=\zeta^{-2}$ and
$\chi_2(2/q)=\chi_2(1/\zeta)$. In case (B), $(-1)^t=(-1)^k\chi_4(s)$ with $\chi_4(s)=(-1)^{(s-1)/2}$, so the sum is
\[
  (-1)^kq^{2k+1}\Bigl(\Ti_2(1/q)-\sum_{s\le2k+1,\ s\text{ odd}}\frac{\chi_4(s)q^{-s}}{s^2}\Bigr)=(-1)^kq^{2k+1}\Ti_2(1/q)-S_k^{\Ti}(q).
\]
All series converge absolutely.
\end{proof}

\subsection{The polynomial part}
For a polynomial $R$, $\mu(R)$ is rational. Indeed $R'(t)=\sum_b\Delta^b[R'](0)\binom tb$ gives the Newton form
\begin{equation}\label{eq:newton}
  \mu(R)=\sum_{b\ge0}\Delta^b[R'](0)\,\Omega_b,\qquad
  \Omega_b:=\sum_{t\ge1}\binom tbw^t=\begin{cases}\dfrac{w}{1-w},&b=0,\\[6pt] \dfrac{w^b}{(1-w)^{b+1}},&b\ge1,\end{cases}
\end{equation}
where $\Delta^b[R'](0)=\sum_{k\le b}(-1)^{b-k}\binom bkR'(k)$ is an integer combination of $R'(0),\dots,R'(b)$, and the sum
is finite ($b<\deg R$). With $y=q$, $iq$ or $q/2$ in cases (A), (B), (C), so that $w=y^{-2}$,
\begin{equation}\label{eq:Omega}
  \Omega_0=\frac1{y^2-1},\qquad \Omega_b=\frac{y^2}{(y^2-1)^{b+1}}\quad(b\ge1).
\end{equation}

\subsection{The pole expansion}

\begin{lemma}[pole expansion]\label{lem:pole}
For $g\in\QQ[u]$ write $g(u)u=D(u)R_g(u)+r_g(u)$ with $\deg r_g<h$. Then $\Lambda(g)=A(g)+\theta\,B(g)$, where
\[
  A(g):=\mu(R_g)+4\sum_{k<h}g(x_k)\,\omega_k\,s_k,\qquad
  B(g):=-4\sum_{k<h}g(x_k)\,\omega_k\,c_k .
\]
Both are rational and $\QQ$-linear in $g$.
\end{lemma}

\begin{proof}
By Lagrange interpolation at the nodes $x_k$,
\[
  r_g(u)=\sum_k\frac{r_g(x_k)\,D(u)}{D'(x_k)\,(u-x_k)},\qquad r_g(x_k)=g(x_k)\,x_k .
\]
Hence $g\varphi=R_g+\sum_kg(x_k)\omega_k/(u-x_k)$, and Lemma~\ref{lem:poleseries} gives the pole terms.
\end{proof}

For an indeterminate $X$ put $\Lambda_X:=A+X\,B$ and
\[
  H(X):=\bigl(\Lambda_X(u^{i+j})\bigr)_{0\le i,j<h},\qquad \Delta_h(X):=\det H(X)\in\QQ[X].
\]
The entries are linear in $X$, so $\deg\Delta_h\le h$, and by Lemma~\ref{lem:pole}
\begin{equation}\label{eq:value}
  \Delta_h(\theta)=\det\bigl(\Lambda(u^{i+j})\bigr)_{0\le i,j<h}.
\end{equation}
When the case matters we write $\Delta^{\chi}_h(X;q)$, $\Delta^{\Ti}_h(X;q)$ and $\Delta^{\chi}_h(X;q/2)$.

\subsection{Change of basis}
If $a_0,\dots,a_{h-1}$ and $b_0,\dots,b_{h-1}$ are polynomials of degree $<h$ with coefficient matrices $C_a$ and $C_b$,
then by linearity of $\Lambda_X$
\begin{equation}\label{eq:gram}
  \bigl(\Lambda_X(a_ib_j)\bigr)_{i,j}=C_a\,H(X)\,C_b^{\mathsf T},\qquad
  \det\bigl(\Lambda_X(a_ib_j)\bigr)=\det C_a\det C_b\,\Delta_h(X).
\end{equation}
We use three bases. The \emph{Lagrange basis} $f_i:=D/(u-x_i)$ has $f_i(x_k)=\ind{i=k}D'(x_i)$; with the Vandermonde
matrix $V=(x_k^d)_{d,k<h}$ we have $C_fV=\operatorname{diag}(D'(x_i))$ and $\prod_iD'(x_i)=\pm(\det V)^2$, where
$\det V=\pm\prod_{k<k'}(k'-k)$. The \emph{anchored basis}
\[
  a_i:=\prod_{m<i}(u-x_m),\qquad b_j:=\prod_{h-j\le m<h}(u-x_m)
\]
consists of monic polynomials of degrees $i$ and $j$, so $\det(\Lambda_X(a_ib_j))=\Delta_h(X)$. The third basis, adapted
to residue classes, is introduced in Section~\ref{sec:classbound}.

% ===========================================================================
\section{The criterion}\label{sec:criterion}
% ===========================================================================

\begin{proposition}[Lemma~E]\label{prop:criterion}
Let $\xi\in\RR$, and for $h\ge1$ let $P_h\in\QQ[X]$ with $\deg P_h\le h$ and let $d_h>0$ be rationals with
$d_hP_h\in\ZZ[X]$. Suppose that there is $c>0$ with $d_h|P_h(\xi)|\le e^{-ch^2}$ for all large $h$, and that for
every $a\in\ZZ$ and $b\ge1$ we have $P_h(a/b)\ne0$ for infinitely many $h$. Then $\xi$ is irrational.
\end{proposition}

\begin{proof}
Suppose $\xi=a/b$. For every $h$, $N_h:=b^hd_hP_h(a/b)$ is an integer. Choose $h$ with $P_h(a/b)\ne0$,
$d_h|P_h(\xi)|\le e^{-ch^2}$ and $b^he^{-ch^2}<1$. Then $1\le|N_h|<1$, a contradiction.
\end{proof}

\begin{theorem}\label{thm:conditional}
Fix a point as in Section~\ref{sec:data}. Suppose there are rationals $d_h>0$ and real numbers $\alpha,\beta$ with
$\alpha+\beta<0$ such that
\begin{enumerate}
\item $d_h\Delta_h\in\ZZ[X]$ for every $h$, and for every $\eps>0$, $\log d_h\le(\alpha+\eps)h^2$ for all large $h$;
\item for every $\eps>0$, $|\Delta_h(\theta)|\le e^{(\beta+\eps)h^2}$ for all large $h$;
\item for every $a\in\ZZ$ and $b\ge1$, $\Delta_h(a/b)\ne0$ for infinitely many $h$.
\end{enumerate}
Then $\theta$ is irrational.
\end{theorem}

\begin{proof}
Apply Proposition~\ref{prop:criterion} with $P_h=\Delta_h$ and $c=-(\alpha+\beta)/2$, using (i) and (ii) with
$\eps=c/2$.
\end{proof}

% ===========================================================================
\section{Denominators I: the symbolic machine and the generic primes}\label{sec:arith1}
% ===========================================================================

\subsection{The free ring}\label{sec:free}
We treat $q$ and $1/(q^2-1)$ as independent indeterminates. Let $\cR:=\QQ[q,E]$. In $\cR$ put $\Omega_0:=E$ and
$\Omega_b:=q^2E^{b+1}$ for $b\ge1$, and for $R\in\QQ[u]$ define $\mu^{\cR}(R):=\sum_b\Delta^b[R'](0)\,\Omega_b\in\cR$ as
in \eqref{eq:newton}. Replacing $\mu$ by $\mu^{\cR}$, $s_k$ by $S_k(q)\in\QQ[q]$ and $c_k$ by $q^{2k+1}$,
Lemma~\ref{lem:pole} defines functionals $A^{\cR},B^{\cR}$ and
\[
  \Delta^{\cR}_h(X):=\det\bigl(A^{\cR}(u^{i+j})+X\,B^{\cR}(u^{i+j})\bigr)_{i,j<h}\in\cR[X].
\]
For a rational $y$ with $y^2\ne1$ let $\ev_y\colon\cR\to\QQ$ be the ring homomorphism with $q\mapsto y$ and
$E\mapsto1/(y^2-1)$. By \eqref{eq:Omega} it maps $\mu^{\cR}$ to $\mu$ and $\Delta^{\cR}_h$ to $\Delta_h$ in case (A)
with $y=q$ and in case (C) with $y=q/2$:
\begin{equation}\label{eq:spec}
  \ev_q(\Delta^{\cR}_h)=\Delta^\chi_h(X;q),\qquad \ev_{q/2}(\Delta^{\cR}_h)=\Delta^\chi_h(X;q/2).
\end{equation}
Case (B) is the specialisation at $y=iq$ followed by $X\mapsto-iX$; Section~\ref{sec:twist} makes this precise over $\QQ$.

\emph{Gauss norms.} For a prime $p$ and $f\in\cR$ let $\|f\|_p$ be the maximum of $|c|_p$ over the rational
coefficients $c$ of $f$, and extend this coefficientwise to $\cR[X]$. Then $\|f+g\|_p\le\max(\|f\|_p,\|g\|_p)$ and
$\|fg\|_p\le\|f\|_p\|g\|_p$. If $|v_0|_p\le1$ and $|v_1|_p\le1$, the evaluation $q\mapsto v_0$, $E\mapsto v_1$
does not increase the norm. By the Leibniz expansion:

\begin{lemma}\label{lem:leibniz}
Let $M$ be an $h\times h$ matrix over $\cR[X]$ (or over $\QQ[X]$) with $\|M_{ij}\|_p\le\rho_i\kappa_j$ for some
$\rho_i,\kappa_j\ge0$. Then $\|\det M\|_p\le\prod_i\rho_i\prod_j\kappa_j$.
\end{lemma}

Two facts carry all the dependence on $q$:
\begin{itemize}
\item[(F1)] $\|\Omega_b\|_p=1$ for all $b$ and $p$, so for $R\in\QQ[u]$ we have
  $\|\mu^{\cR}(R)\|_p\le\sup_{t\in\ZZ_{\ge0}}|R'(t)|_p$;
\item[(F2)] $\|q^{2k+1}X-S_k(q)\|_p\le p^{2\#\{e\ge1:\,p^e\le2k+1\}}$, since the coefficients of $S_k$ are $1/s^2$
  with $s$ odd, $s\le2k+1$, and $v_p(s)\le\#\{e\ge1:p^e\le s\}$.
\end{itemize}

\emph{The Gram identity.} Write $\Lambda^{\cR}_X:=A^{\cR}+XB^{\cR}$. By the argument of \eqref{eq:gram}, for any family
$(e_m)_{m<h}$ of polynomials of degree $<h$ with coefficient matrix $C_e$,
\begin{equation}\label{eq:gramR}
  \det\bigl(\Lambda^{\cR}_X(e_me_{m'})\bigr)_{m,m'<h}=(\det C_e)^2\,\Delta^{\cR}_h(X).
\end{equation}
For a pole $x_l$ we have $\Lambda^{\cR}_X(g)=\mu^{\cR}(R)-4c\,(q^{2l+1}X-S_l(q))$ whenever $g\varphi=R+c/(u-x_l)$ with $R$
a polynomial; more generally the pole parts add up.

Below, $\lev(p):=\#\{e\ge1:p^e\le2h-1\}$. For $t\in\ZZ_{\ge0}$ and $m<h$ we have $t-x_m=(2t+2m+1)/2$, and for $p$ odd
$v_p(t-x_m)=v_p(2t+2m+1)$. The numbers $2t+2m+1$, $m<h$, are $h$ consecutive odd integers, so for each level
$d=p^e$ those divisible by $d$ are the $m$ in one residue class modulo $d$.

\subsection{All odd primes: the Lagrange basis}

\begin{lemma}[Lemma~9]\label{lem:nine}
For every odd prime $p$ and every $h\ge1$, $\|\Delta^{\cR}_h\|_p\le p^{2h\lev(p)}$.
\end{lemma}

\begin{proof}
Use the Lagrange basis $f_i=D/(u-x_i)$ and put $G_{ij}:=\Lambda^{\cR}_X(f_if_j)$. Since $C_fV=\operatorname{diag}(D'(x_i))$
and $\prod_iD'(x_i)=\pm(\det V)^2$, we get $(\det C_f)^2=\pm\prod_iD'(x_i)$, so by \eqref{eq:gramR}
\[
  \|\Delta^{\cR}_h\|_p=\|\det G\|_p\Big/\prod_i|D'(x_i)|_p .
\]
By Lemma~\ref{lem:leibniz} it suffices to prove the \emph{row claim} $\|G_{ij}\|_p\le|D'(x_i)|_p\,p^{2\lev(p)}$. For a
level $e\ge1$ put $N_e(i):=\#\{m<h:m\ne i,\ p^e\mid m-i\}$; then $v_p(D'(x_i))=\sum_eN_e(i)$, and $N_e(i)=0$ unless
$p^e\le h-1$. So it suffices to show that every value $R'(t)$ entering (F1) has
\begin{equation}\label{eq:rowclaim}
  v_p(R'(t))\ge\sum_e\bigl(N_e(i)-2\ind{N_e(i)\ge1}\bigr),
\end{equation}
and that the pole part obeys the same bound.

Fix $t\in\ZZ_{\ge0}$ and a level $d=p^e$, and let $m_0$ be the residue class mod $d$ of the indices $m$ with
$d\mid2t+2m+1$. The class of $i$ has $N_e(i)+1$ elements in $[0,h)$, and residue classes mod $d$ in $[0,h)$ differ in
size by at most one, so the class $m_0$ has at least $N_e(i)$ elements.

\emph{Off-diagonal entries} ($i\ne j$). Here $f_if_j\varphi=F_{ij}:=u\prod_{m\ne i,j}(u-x_m)$ is a polynomial, so
$G_{ij}=\mu^{\cR}(F_{ij})$. Each term of $F'_{ij}(t)$ is a product of all factors but one. At level $d$ the factors
$t-x_m$ divisible by $d$ are those with $m\equiv m_0$, $m\ne i,j$; omitting one factor leaves at least
$N_e(i)-2$ of them in both cases $m_0\equiv i$ and $m_0\not\equiv i$. Summing over levels ($v_p$ of a product is
$\sum_e$ of the number of factors divisible by $p^e$) gives \eqref{eq:rowclaim}.

\emph{Diagonal entries.} Put $g_i:=u\prod_{m\ne i}(u-x_m)$. Then $f_i^2\varphi=g_i/(u-x_i)=Q_i+g_i(x_i)/(u-x_i)$ with
$Q_i:=(g_i(u)-g_i(x_i))/(u-x_i)$ and $g_i(x_i)=x_iD'(x_i)$. The pole part is $-4x_iD'(x_i)(q^{2i+1}X-S_i)$; by (F2),
$|x_i|_p\le1$ and $2i+1\le2h-1$, its norm is at most $|D'(x_i)|_p\,p^{2\lev(p)}$. For the polynomial part, list the
$h$ roots $\{x_m:m\ne i\}\cup\{0\}$ of $g_i$ as $\alpha_1,\dots,\alpha_h$; telescoping
$(u-\alpha_r)-(x_i-\alpha_r)=u-x_i$ gives
\[
  Q_i(u)=\sum_{r=1}^{h}\prod_{s<r}(x_i-\alpha_s)\prod_{s>r}(u-\alpha_s).
\]
Order the roots by decreasing $v_p(m-i)$, with $0$ last. Then for every level $e$ simultaneously, the $\alpha_s=x_m$ with
$p^e\mid m-i$ form an initial segment of length $N_e(i)$ (note $x_i-x_m=m-i$). Fix $t$, a level $d=p^e$ and a term
$r$ of $Q_i'(t)$ (a product with one more factor omitted). If $r>N_e(i)$, the first product already has $N_e(i)$
factors divisible by $d$. If $r\le N_e(i)$, it has $r-1$; if moreover $m_0\equiv i$, the remaining $N_e(i)-r$ class
members sit at positions $s>r$ and at most one of them is omitted, giving at least $N_e(i)-2$ in total; if
$m_0\not\equiv i$, the whole class $m_0$ (at least $N_e(i)$ members, all with $v_p(m-i)<e$) sits at positions
$s>N_e(i)\ge r$, giving at least $N_e(i)-1$. So \eqref{eq:rowclaim} holds for $Q_i'(t)$ as well. Finally
$\#\{e:N_e(i)\ge1\}\le\lev(p)$.
\end{proof}

\subsection{One-level primes: a basis adapted to residue classes}\label{sec:classbound}
For a modulus $d$ and $m<h$ write $r_m:=\lfloor m/d\rfloor$ (the rank of $m$ in its class), $n_m$ for the size of the
class of $m$ modulo $d$ in $[0,h)$, and put
\begin{gather*}
  w_d(l):=2\ind{d\le2l+1}-\ind{d\mid2l+1},\\
  M_m:=\max\{w_d(l):l\le m,\ l\equiv m\ (d)\},\\
  \cb(h,d):=\sum_{m<h}\bigl(2r_m+1+M_m-n_m\bigr)_+ .
\end{gather*}

\begin{proposition}[class bound]\label{prop:classbound}
Let $p$ be an odd prime with $p^2>2h$. Then $\|\Delta^{\cR}_h\|_p\le p^{\cb(h,p)}$.
\end{proposition}

\begin{proof}
Since $p^2>2h$, every difference $m-l$ of distinct indices and every odd $s\le2h-1$ has $v_p\le1$; so
$v_p(m-l)=\ind{p\mid m-l}$ and $v_p(2l+1)=\ind{p\mid2l+1}$. For $m<h$ let $\operatorname{seg}(m):=\{l\le m:l\equiv m\ (p)\}$
(it has $r_m+1$ elements), and put
\[
  e_m:=s_m\,\frac{D(u)}{\prod_{l\in\operatorname{seg}(m)}(u-x_l)},\qquad s_m:=\prod_{l\in\operatorname{seg}(m),\,l<m}(x_m-x_l).
\]
Then $\deg e_m=h-r_m-1<h$ and $v_p(s_m)=r_m$. The polynomial $e_m$ vanishes at $x_k$ unless $k\in\operatorname{seg}(m)$, and
$e_m(x_m)=D'(x_m)$. So the evaluation matrix $(e_m(x_k))=C_eV$ is lower triangular with diagonal $D'(x_m)$, hence
$\det C_e\det V=\prod_mD'(x_m)=\pm(\det V)^2$ and $\det C_e=\pm\det V$. By \eqref{eq:gramR},
$\|\Delta^{\cR}_h\|_p=\|\det G^e\|_p/|\det V|_p^2$ with $G^e_{mm'}:=\Lambda^{\cR}_X(e_me_{m'})$, and
$2v_p(\det V)=2\#\{k<k':p\mid k'-k\}=2\sum_mr_m$.

Let $\cA:=\operatorname{seg}(m)$, $\mathcal B:=\operatorname{seg}(m')$, $r=r_m$, $r'=r_{m'}$. Then
\[
  e_me_{m'}\varphi=s_ms_{m'}\,u\,\frac{\prod_{l\notin\cA\cup\mathcal B}(u-x_l)}{\prod_{l\in\cA\cap\mathcal B}(u-x_l)} .
\]
If $m$ and $m'$ lie in different classes, $\cA\cap\mathcal B=\emptyset$, the right side is $s_ms_{m'}$ times a polynomial in
$\ZZ[\frac12][u]$, and (F1) gives $v_p(G^e_{mm'})\ge r+r'$. If they lie in the same class, of size $n$, say with
$r\le r'$, then $\cA\subseteq\mathcal B$. Division by the monic polynomial $\prod_{l\in\cA}(u-x_l)\in\ZZ[\frac12][u]$ leaves a
$p$-integral polynomial part, again of valuation $\ge r+r'$. The residue at $x_l$, $l\in\cA$, is
\begin{gather*}
  c_l=s_ms_{m'}\,x_l\,\frac{\prod_{k\notin\mathcal B}(x_l-x_k)}{\prod_{k\in\cA,\,k\ne l}(x_l-x_k)},\\
  v_p(c_l)=r+r'+\ind{p\mid2l+1}+(n-1-r')-r,
\end{gather*}
because exactly the $n-1-r'$ class members after $m'$ lie outside $\mathcal B$, and the $r$ other elements of $\cA$ lie in
the class. By (F2) and the single level, the pole factor $q^{2l+1}X-S_l$ has norm at most $p^{2\ind{p\le2l+1}}$. Hence
\[
  v_p(G^e_{mm'})\ge r+r'+\beta(r,r'),\qquad \beta(a,b):=\min\bigl(0,\ n-1-a-b-\widehat M_{\min(a,b)}\bigr),
\]
where $\widehat M_a$ is the largest weight $w_p(l)$ among the class members of rank $\le a$. Since $\widehat M$ is
nondecreasing and $x\mapsto\min(0,x)$ is concave, $\beta(a,b)\ge\frac12(\beta(a,a)+\beta(b,b))$. So every entry, in
the same class or not, satisfies $v_p(G^e_{mm'})\ge\rho_m+\rho_{m'}$ with $\rho_m:=r_m+\frac12\beta(r_m,r_m)\le r_m$.
Lemma~\ref{lem:leibniz} gives $v_p(\det G^e)\ge2\sum_m\rho_m$, and therefore
\[
  v_p(\Delta^{\cR}_h)\ge\sum_m\beta(r_m,r_m)=-\sum_m(2r_m+1+M_m-n_m)_+=-\cb(h,p).\qedhere
\]
\end{proof}

\subsection{The class count}
Let $\varphi_1$ be the continuous, piecewise linear, nonincreasing function on $[2/5,\infty)$ with
\[
  \varphi_1(x)=\begin{cases}3-\tfrac52x,&\tfrac25\le x\le\tfrac12,\\ 2-\tfrac12x,&\tfrac12\le x\le\tfrac23,\\
  3-2x,&\tfrac23\le x\le1,\\ 2-x,&1\le x\le2,\end{cases}
\]
and $\varphi_1(x)=0$ for $x\ge2$. Then $\int_{2/5}^{2}\varphi_1=\frac3{16}+\frac{41}{144}+\frac49+\frac12=\frac{17}{12}$.

\begin{lemma}[Lemma~8A, class count]\label{lem:classcount}
Let $h\ge2$ and let $d$ be odd. If $2h\le5d$ and $d<2h$ then $\cb(h,d)\le h\,\varphi_1(d/h)$; more precisely
$\cb(h,d)=h\varphi_1(d/h)-\frac32$ if $d\le(2h-1)/3$ and $\cb(h,d)=h\varphi_1(d/h)$ otherwise. If $d\ge2h$ then $\cb(h,d)=0$.
(The formal version assumes $h\ge13$, which is needed only in Theorem~\ref{thm:generic}.)
\end{lemma}

\begin{proof}
Write $K=h$ and $\kappa=(d-1)/2$. If $d\ge2K$, every atom $l<K$ has $2l+1<d$, so all weights vanish, every class is a
singleton and every term is $(1+0-1)_+=0$. Now let $2K\le5d$ and $d<2K$; as $d$ is odd, $5d\ge2K+1$ and $3d\ne2K$. The
odd multiples of $d$ below $2K$ are $d$ and possibly $3d$. Hence an atom $l$ has weight $0$ if $l<\kappa$ (we call it
\emph{small}), weight $1$ if $l=\kappa$ or $l=\kappa+d$, and weight $2$ otherwise. Every atom after the first in its
class is $\ge d>\kappa$, so only the first atom $c$ of a class can be small. We evaluate $\cb$ class by class; a class
$\{c,c+d,\dots\}$ of size $n$ contributes $\sum_r(2r+1+M-n)_+$ over its atoms.
\begin{itemize}
\item[(IV)] $K<d<2K$. All classes are singletons, contributing their weights: $2(K-\kappa)-1=2K-d$.
\item[(III)] $(2K+1)/3\le d\le K$. Classes $c<K-d$ are pairs, with $c<\kappa$ small; they contribute $0+(1+2)=3$.
  The other classes $c\in[K-d,d-1]$ are singletons; those with $c\ge\kappa$ ($(d+1)/2$ of them, one of weight $1$)
  contribute $d$ in total. The atom $\kappa+d$ is $\ge K$. Total $3(K-d)+d=3K-2d$.
\item[(II)] $K/2\le d\le(2K-1)/3$. Classes $c<K-d$ are pairs and $c\in[K-d,d-1]$ singletons, all of weight $2$ (as
  $\kappa<K-d$). A pair contributes $(w(c)-1)_++1+M$: this is $3$ for $c<\kappa$, $2$ for $c=\kappa$ (weights
  $1,1$), and $4$ for $\kappa<c<K-d$. Total $3\kappa+2+4(K-d-\kappa-1)+2(2d-K)=2K-\frac d2-\frac32$.
\item[(I)] $(2K+1)/5\le d\le K/2$. Classes $c<K-2d$ are triples, with $c<\kappa$ (as $K-2d\le\kappa$), contributing
  $0+2+4=6$. Classes $c\in[K-2d,d-1]$ are pairs, contributing $3$ for $c<\kappa$ ($\frac{5d-1}2-K$ of them), $2$ for
  $c=\kappa$, and $4$ for $c>\kappa$ ($\frac{d-1}2$ of them). Total $3K-\frac52d-\frac32$.
\end{itemize}
In the four regions $K\varphi_1(d/K)$ equals $2K-d$, $3K-2d$, $2K-\frac d2$ and $3K-\frac52d$.
\end{proof}

\subsection{The generic primes}
For odd $p$ put
\[
  e(h,p):=\begin{cases}\cb(h,p),& p^2>2h\text{ and }5p\ge2h,\\ 2h\lev(p),&\text{otherwise.}\end{cases}
\]

\begin{theorem}[generic primes]\label{thm:generic}
For every odd prime $p$ and every $h$, $\|\Delta^{\cR}_h\|_p\le p^{e(h,p)}$. For every $\eps>0$ and all large $h$,
\[
  \sum_{p\text{ odd}}e(h,p)\log p\le\Bigl(\frac{133}{60}+\eps\Bigr)h^2 .
\]
\end{theorem}

\begin{proof}
The first statement is Proposition~\ref{prop:classbound} or Lemma~\ref{lem:nine}. For the sum let $h\ge13$. If
$5p\ge2h$ then $p^2\ge4h^2/25>2h$, so $e(h,p)=\cb(h,p)$, which vanishes for $p\ge2h$ and is at most $h\varphi_1(p/h)$
for $2h/5\le p<2h$ (Lemma~\ref{lem:classcount}). So the sum is at most
\[
  \sum_{p<2h/5}2h\lev(p)\log p+\sum_{2h/5\le p<2h}h\,\varphi_1(p/h)\log p .
\]
For the first sum, $\lev(p)\le1$ when $p^2>2h-1$ and $\lev(p)\log p\le\log(2h)$ always, so it is at most
$2h\bigl(\theta(2h/5)+(\sqrt{2h}+1)\log(2h)\bigr)=(\frac45+o(1))h^2$ by the prime number theorem $\theta(x)\sim x$.
For the second, cut $[2/5,2]$ into $M$ equal cells. On each cell $\varphi_1$ is at most its value at the left end,
and $\sum_{p\in hI}\log p=h|I|+o(h)$ for each fixed cell $I$. Hence the second sum is at most
$h^2\bigl(\int_{2/5}^2\varphi_1+\frac{16}{5M}\bigr)+o(h^2)$, and $M$ is arbitrary.
\end{proof}

\subsection{The twist}\label{sec:twist}
For case (B) let $\ev^{\Ti}_{q}\colon\cR\to\QQ$ be the homomorphism with $q\mapsto q$ and $E\mapsto1/(q^2+1)$, for an
integer $q\ge2$.

\begin{proposition}[the twist]\label{prop:twist}
\begin{enumerate}
\item Every monomial $q^{m_0}E^{m_1}X^j$ occurring in $\Delta^{\cR}_h$ has $m_0+j$ even.
\item Let $\Delta^{\cR,\Ti}_h\in\cR[X]$ be obtained from $\Delta^{\cR}_h$ by multiplying the coefficient of each monomial
  $q^{m_0}E^{m_1}X^j$ by $\tau(m,j):=(-1)^{m_1+j+(m_0+j)/2}$. Then $\ev^{\Ti}_q(\Delta^{\cR,\Ti}_h)=\Delta^{\Ti}_h(X;q)$.
\end{enumerate}
\end{proposition}

\begin{proof}
(i) The automorphism $q\mapsto-q$, $X\mapsto-X$ of $\cR[X]$ fixes every entry $A^{\cR}(u^n)+XB^{\cR}(u^n)$: $\Omega_b$ and
$S_k(q)$ are even in $q$, and $q^{2k+1}$ is odd. So it fixes $\Delta^{\cR}_h$.

(ii) Over $\CC$, the homomorphism $q\mapsto iq$, $E\mapsto-1/(q^2+1)=1/((iq)^2-1)$ maps $\Omega_b$ to the values
\eqref{eq:Omega} with $y=iq$, i.e.\ to those of case (B); it maps $S_k(q)$ to $S_k(iq)=S^{\Ti}_k(q)$, and $q^{2k+1}$ to
$i(-1)^kq^{2k+1}$. So it maps $A^{\cR}$ to the functional $A$ of case (B) and $B^{\cR}$ to $i$ times its $B$, hence
$\Delta^{\cR}_h(X)$ to $\Delta^{\Ti}_h(iX;q)$. Replacing $X$ by $-iX$, the monomial $c\,q^{m_0}E^{m_1}X^j$ goes to
$c\,i^{m_0}(-1)^{m_1}(-i)^jq^{m_0}(q^2+1)^{-m_1}X^j$, and by (i), $i^{m_0}(-i)^j(-1)^{m_1}=\tau(m,j)$.
\end{proof}

So the coefficients of $\Delta^{\cR,\Ti}_h$ are, up to sign, those of $\Delta^{\cR}_h$: the two have the same Gauss
norms, the same $E$-degrees and the same $q$-degrees, coefficient by coefficient.

% ===========================================================================
\section{\texorpdfstring{Denominators II: the odd primes dividing $q$ and the pole factor}{Denominators II: the odd primes dividing q and the pole factor}}\label{sec:arith2}
% ===========================================================================

\subsection{The pole order}
For $f\in\cR$ let $\Ed f$ be its degree in $E$. Put
\begin{gather*}
  \rho_i:=(i+1-\lceil h/2\rceil)_+,\qquad \kappa_j:=(j-\lfloor h/2\rfloor)_+,\\
  \beta_h:=\sum_{i<h}\rho_i+\sum_{j<h}\kappa_j=\lfloor h^2/4\rfloor\le\frac{h^2+2h}4 .
\end{gather*}

\begin{lemma}[Lemma~5]\label{lem:poleorder}
Every coefficient of $\Delta^{\cR}_h\in\cR[X]$ has $E$-degree at most $\beta_h$.
\end{lemma}

\begin{proof}
In the monomial basis, the entry $(i,j)$ is $\mu^{\cR}(R)$ plus pole terms that do not involve $E$, where $R$ is the
quotient of $u^{i+j+1}$ by $D$, of degree $i+j+1-h$ (and $R=0$ if $i+j+1<h$). In $\mu^{\cR}(R)=\sum_{b<\deg R}\Delta^b[R'](0)\Omega_b$
each $\Omega_b$ has $E$-degree $b+1\le\deg R$. So the entry has $E$-degree at most $(i+j+1-h)_+\le\rho_i+\kappa_j$, and
the Leibniz expansion gives the claim.
\end{proof}

\emph{Clearing $E$.} For $f=\sum c_{a,b}q^aE^b\in\cR$ with $\Ed f\le\beta$ and $\sigma\in\{\pm1\}$ put
$\cl^\sigma_\beta(f):=\sum c_{a,b}\,q^a(q^2-\sigma)^{\beta-b}\in\QQ[q]$, and $\cl_\beta:=\cl^{+}_\beta$. Since
$(q^2\mp1)^k$ has integer coefficients, $\|\cl^\sigma_\beta f\|_p\le\|f\|_p$ for every $p$. If $v_1(v_0^2-\sigma)=1$,
then evaluating at $q=v_0$, $E=v_1$ gives
\begin{equation}\label{eq:clear}
  \ev\bigl(\cl^\sigma_\beta f\bigr)=(v_0^2-\sigma)^\beta\,\ev(f).
\end{equation}
By the computation in the proof of Proposition~\ref{prop:twist}(ii), for $f=[X^j]\Delta^{\cR}_h$ and
$f^{\Ti}=[X^j]\Delta^{\cR,\Ti}_h$, the coefficient of $q^n$ in $\cl^{-}_{\beta}(f^{\Ti})$ is $\pm$ the coefficient of
$q^n$ in $\cl_\beta(f)$: indeed $\cl^{-}_{\beta}(f^{\Ti})(q)$ is a unit of $\{\pm1,\pm i\}$ times $\cl_\beta(f)(iq)$.

\subsection{The \texorpdfstring{$q$}{q}-adic gain}
Let $\psi_q\colon\cR\to\QQ[[q]]$ be the ring homomorphism with $q\mapsto q$ and $E\mapsto1/(q^2-1)=-\sum_{n\ge0}q^{2n}$. Put
\[
  a_h:=\max\Bigl(0,\Bigl\lfloor\frac{h(h-1)-6h}{3}\Bigr\rfloor\Bigr)\ge\frac{h^2-7h-2}{3} .
\]

\begin{lemma}[Lemma~4$'$]\label{lem:qadic}
Every coefficient of $\psi_q(\Delta^{\cR}_h)\in\QQ[[q]][X]$ is divisible by $q^{a_h}$. Consequently, for every $k$, all
monomials of $\cl_{\beta_h}\bigl([X^k]\Delta^{\cR}_h\bigr)\in\QQ[q]$ have degree at least $a_h$.
\end{lemma}

\begin{proof}
\emph{Reflection.} From $\psi_q(\Omega_b)=(-1)^{b+1}q^2(1-q^2)^{-b-1}$ for $b\ge1$ and $\psi_q(\Omega_0)=-(1-q^2)^{-1}$, the
coefficient of $q^{2m}$ in $\psi_q(\Omega_b)$ is $-(-1)^b\binom{m+b-1}{b}$, and odd powers do not occur. Newton's backward
formula $f(-m)=\sum_b(-1)^b\binom{m+b-1}b\Delta^bf(0)$ for polynomials $f$ then gives
\begin{equation}\label{eq:reflect}
  \psi_q\bigl(\mu^{\cR}(R)\bigr)=-\sum_{m\ge0}q^{2m}R'(-m)\qquad(R\in\QQ[u]).
\end{equation}
\emph{Recombination.} For a polynomial $g$ write $g\varphi=R+\sum_kc_k'/(u-x_k)$, $c_k'=g(x_k)\omega_k$. Since
$\bigl((u-x_k)^{-1}\bigr)'(-m)=-4/(2k+1-2m)^2$ and
\[
  \sum_{m\ge0}\frac{q^{2m}}{(2k+1-2m)^2}=S_k(q)+q^{2k+1}\chi_2(q),
\]
where $\chi_2(q)=\sum_{j\ge0}q^{2j+1}/(2j+1)^2\in\QQ[[q]]$, \eqref{eq:reflect} and Lemma~\ref{lem:pole} combine to
\begin{equation}\label{eq:recomb}
  \psi_q\bigl(\Lambda^{\cR}_X(g)\bigr)=-\sum_{m\ge0}q^{2m}(g\varphi)'(-m)-4\,(X+\chi_2(q))\sum_{k<h}g(x_k)\,\omega_k\,q^{2k+1}.
\end{equation}
\emph{Slots.} For $g=u^{i+j}$ we have
$(g\varphi)'(-m)=\varphi'(-m)\,t^it^j+\varphi(-m)\bigl(it^{i-1}\bigr)t^j+\varphi(-m)\,t^i\bigl(jt^{j-1}\bigr)$ at $t=-m$, and
$g(x_k)=x_k^ix_k^j$. So the matrix $\psi_q(H)$ is a sum of rank-one matrices $c\,xy^{\mathsf T}$, where the row vector $x$ is
either a \emph{value row} $(t^i)_i$ or a \emph{derivative row} $(it^{i-1})_i$ at a node $t=-m$, or a pole row
$(x_k^i)_i$; the same holds for $y$. A node $-m$ carries the factor $q^{2m}$ and a pole $k$ carries $q^{2k+1}$. Give these
row types the \emph{keys} $(m,\mathrm{val})$, $(m,\mathrm{der})$ and $k$, with \emph{prices} $2m$, $2m$ and $2k+1$.

\emph{Truncation and Andreief.} Let $M:=a_h+1$ and let $H_M$ be the same sum restricted to the nodes $m<M$. Then
$\psi_q(H)\equiv H_M$ entrywise modulo $q^{2M}$, so $\psi_q(\Delta^{\cR}_h)=\det\psi_q(H)\equiv\det H_M\pmod{q^{2M}}$. For a finite
sum of rank-one matrices, the Andreief (Cauchy--Binet) identity \cite{Andreief,Forrester2019} reads
\begin{equation}\label{eq:andreief}
  h!\,\det\Bigl(\sum_\iota c_\iota x_\iota y_\iota^{\mathsf T}\Bigr)
  =\sum_{f\colon[h]\to I}\Bigl(\prod_ic_{f(i)}\Bigr)\det\bigl(x_{f(i)}(j)\bigr)_{i,j}\det\bigl(y_{f(i)}(j)\bigr)_{i,j}.
\end{equation}
Here the rank-one terms are indexed by $\iota\in I$; the row vector $x_\iota$ and the power of $q$ in $c_\iota$ depend only
on the key of $\iota$ (the two terms $\varphi'(-m)\,t^it^j$ and $\varphi(-m)\,t^i(jt^{j-1})$ of a node share the key
$(m,\mathrm{val})$). A term of \eqref{eq:andreief} vanishes unless the rows $x_{f(i)}$ are pairwise distinct, so unless
the $h$ keys of $f(0),\dots,f(h-1)$ are pairwise distinct. The term is then divisible by $q$ to the sum of their prices.
The map $(m,\mathrm{val})\mapsto3m$, $(m,\mathrm{der})\mapsto3m+2$, $k\mapsto3k+1$ is injective and satisfies
$2\cdot(\text{image})\le3\cdot\text{price}+6$. The images of $h$ distinct keys are $h$ distinct naturals, with sum
$\ge h(h-1)/2$; hence the prices sum to at least $(h(h-1)-6h)/3\ge a_h$. Since $h!$ is invertible and $2M>a_h$,
$q^{a_h}$ divides every coefficient of $\psi_q(\Delta^{\cR}_h)$.

\emph{Clearing.} Let $f=[X^k]\Delta^{\cR}_h$. By Lemma~\ref{lem:poleorder} and \eqref{eq:clear} (with $v=(q,1/(q^2-1))$ in
$\QQ[[q]]$), $\cl_{\beta_h}f=(q^2-1)^{\beta_h}\psi_q(f)$ in $\QQ[[q]]$, which is divisible by $q^{a_h}$. As
$\cl_{\beta_h}f$ is a polynomial in $q$, its monomials of degree $<a_h$ vanish.
\end{proof}

\subsection{The odd part of the denominator}
For the three cases put
\[
  (y,Y)=(q,\,q^2-1)\ \text{(A)},\qquad (q,\,q^2+1)\ \text{(B)},\qquad (q/2,\,q^2-4)\ \text{(C)},
\]
so that $Y$ is $y^2-1$, $-(y^2-1)$ with $y=iq$, or $4(y^2-1)$, and put
\[
  \alpha_{\odd}:=\frac{133}{60}+\frac14\log\odd(Y)-\frac13\log\odd(q).
\]

\begin{theorem}[the odd primes]\label{thm:oddprimes}
Fix a point of case (A) or (B) with $q\ge3$, or of case (C) with $q\ge5$ odd. There are positive rationals $d^{\odd}_h$,
with $v_2(d^{\odd}_h)=0$, such that $|d^{\odd}_hc|_p\le1$ for every coefficient $c$ of $\Delta_h$ and every odd prime $p$,
and for every $\eps>0$ and all large $h$,
\[
  \log d^{\odd}_h\le(\alpha_{\odd}+\eps)h^2 .
\]
Consequently, if $n_2(h)\in\ZZ$ satisfy $|c|_2\le2^{n_2(h)}$ for every coefficient $c$ of $\Delta_h$, then
$d_h:=2^{n_2(h)}d^{\odd}_h$ satisfies $d_h\Delta_h\in\ZZ[X]$ and $\log d_h=n_2(h)\log2+\log d^{\odd}_h$.
\end{theorem}

\begin{proof}
Let $c=[X^k]\Delta_h$. In cases (A) and (C), $c=\ev_y(f)$ with $f=[X^k]\Delta^{\cR}_h$ by \eqref{eq:spec}; in case (B),
$c=\ev^{\Ti}_q(f^{\Ti})$ with $f^{\Ti}=[X^k]\Delta^{\cR,\Ti}_h$ by Proposition~\ref{prop:twist}. Let $p$ be odd. Then $y$
is $p$-integral, with $v_p(y)=v_p(q)$, and $v_p$ of the value of $E$ is $-v_p(Y)$. We use the clearing $\cl^\sigma$ with
$\sigma=+1$ in cases (A), (C) and $\sigma=-1$ in case (B), so that $\ev(\cl^\sigma f)=(y^2-\sigma)^{\beta_h}c$ and
$v_p(y^2-\sigma)=v_p(Y)$. By Proposition~\ref{prop:twist} and the remark after \eqref{eq:clear}, the Gauss norms and the
$q$-degrees of the monomials of the cleared coefficient are the same as for $\cl_{\beta_h}(f)$ in case (A).
\begin{itemize}
\item $p\mid q$: then $p\nmid Y$, and by Lemma~\ref{lem:qadic} every monomial of the cleared coefficient has degree
  $\ge a_h$, so $|c|_p\le\|f\|_p\,|q|_p^{a_h}\le p^{2h\lev(p)}|q|_p^{a_h}$ by Lemma~\ref{lem:nine}.
\item $p\mid Y$: by Lemma~\ref{lem:poleorder} and \eqref{eq:clear},
  $|c|_p\le\|f\|_p\,|Y|_p^{-\beta_h}\le p^{2h\lev(p)+v_p(Y)\beta_h}$, by Lemma~\ref{lem:nine}.
\item $p\nmid qY$: $y$ and the value of $E$ are $p$-integral, so $|c|_p\le\|f\|_p\le p^{e(h,p)}$ by
  Theorem~\ref{thm:generic}.
\end{itemize}
Let $N_h:=\prod_pp^{e_p}$ over the odd primes, where $e_p=2h\lev(p)+e(h,p)+v_p(Y)\beta_h$ for $p\mid qY$ and $e_p=e(h,p)$
otherwise. Only primes $p<2h+q^2$ occur. Put $d^{\odd}_h:=N_h/\odd(q)^{a_h}$; since $|\odd(q)|_p=|q|_p$ for odd $p$,
$|d^{\odd}_hc|_p\le1$. Finally the terms $2h\lev(p)\log p\le2h\log(2h)$ of the finitely many $p\mid qY$ are $O(h\log h)$,
$\beta_h\log\odd(Y)\le\frac{h^2+2h}4\log\odd(Y)$, $a_h\ge(h^2-7h-2)/3$, and Theorem~\ref{thm:generic} bounds
$\sum_pe(h,p)\log p$.
\end{proof}

\begin{remark}
In \cite{Chi2Paper} the gain at $p=q$ used that $q$ is prime. Here it is applied at each odd prime $p\mid q$ separately,
with multiplicity $v_p(q)$; the $2$-adic share of the gain, for even $q$, is counted inside the $2$-adic lemmas of
Section~\ref{sec:A}, and only there.
\end{remark}

% ===========================================================================
\section{\texorpdfstring{The size of $\Delta_h(\theta)$}{The size of Delta_h(theta)}}\label{sec:analytic}
% ===========================================================================

\subsection{Dual certificates}
Let $\psi(u):=u\log|u|-u$ (with $\psi(0)=0$), an antiderivative of $\log|u|$, so that
$\int_\alpha^\beta\log|x-y|\,dy=\psi(x-\alpha)-\psi(x-\beta)$. Put
for $0<W<1$ and $x\ge0$
\[
  V_W(x):=x\log W-\psi(x+1)+\psi(x)=x\log W-(x+1)\log(x+1)+x\log x+1 .
\]
We write $V=V_W$ when $W$ is fixed.
A \emph{step density} is $\mu=\sum_{j<J}c_j\mathbf 1_{[a_j,a_{j+1})}$ with $0=a_0\le a_1\le\dots\le a_J$, $0\le c_j\le2$ and
$\int\mu=1$. Its potential and energy are
\[
  U^\mu(x):=\int\log|x-y|\,\mu(y)\,dy,\qquad E(\mu):=\iint\log|x-y|\,\mu(x)\mu(y)\,dx\,dy .
\]

\begin{definition}\label{def:cert}
A \emph{dual certificate} for a weight $W\in(0,1)$ consists of a step density $\mu$, a real number $\lambda$, integers $L,X,N\ge1$ with
$a_J\le N\le X$, and real numbers $S_k$ ($0\le k<XL$) such that
\begin{enumerate}
\item $F(x):=2U^\mu(x)+V(x)\le S_k$ for $x\in[k/L,(k+1)/L]$, for every $k<XL$;
\item $1/X+\log W<0$ and $\log X+X\log W<\lambda$.
\end{enumerate}
Its \emph{exponent} is $\Phi:=-E(\mu)+\lambda+\frac2L\sum_{k<XL}(S_k-\lambda)_+$.
\end{definition}

\begin{theorem}[the size bound]\label{thm:size}
Fix a point of case (A), (B) or (C) and put $W:=|w|$. For every dual certificate for $W$, as $h\to\infty$,
\[
  \log|\Delta_h(\theta)|\le\Phi h^2+O(h\log h).
\]
\end{theorem}

The constant in $O(h\log h)$ depends on the certificate. The heuristic behind the theorem: by \eqref{eq:value},
$\Delta_h(\theta)$ is a Hankel determinant of a functional on the nodes $t=1,2,\dots$ with the
derivative structure of $(g\varphi)'$ and weights of size $W^t$, i.e.\ a discrete log-gas in the external field $V$ with at most two
particles per node. The dual certificate is a feasible point of the dual of the corresponding constrained
equilibrium problem \cite[Ch.~I]{SaffTotik}, and the proof below is a discrete version of weak duality.

\subsection{Node expansion}
We have $\Lambda(u^n)=\sum_{s\ge1}w^s\bigl(ns^{n-1}\varphi(s)+s^n\varphi'(s)\bigr)$. Let $H^{(T)}$ be the matrix of
truncated sums over $1\le s\le T$. Then $H^{(T)}\to(\Lambda(u^{i+j}))_{i,j}$ entrywise, so by continuity of the
determinant $|\Delta_h(\theta)|\le\sup_T|\det H^{(T)}|$. Expanding $(i+j)s^{i+j-1}=(is^{i-1})s^j+s^i(js^{j-1})$,
\[
  H^{(T)}=\sum_{s\le T}\Bigl(w^s\varphi'(s)\,v_sv_s^{\mathsf T}+w^s\varphi(s)\,v'_sv_s^{\mathsf T}+w^s\varphi(s)\,v_sv_s'^{\mathsf T}\Bigr),
\]
with the value row $v_s=(s^j)_{j<h}$ and the derivative row $v'_s=(js^{j-1})_{j<h}$. We call $(s,\mathrm{val})$ and
$(s,\mathrm{der})$ the two \emph{slots} of the node $s$. Since $\varphi'/\varphi=1/s-\sum_k1/(s-x_k)$ is a difference of two
numbers in $[0,h]$ for $s\ge1$, all three coefficients are bounded by
\[
  b(s):=W^s\,h\,\varphi(s),
\]
since $|w^s|=W^s$.

\begin{lemma}[minors with derivative rows]\label{lem:slotdet}
Let $h\ge2$ and let $(s_i,\tau_i)_{i<h}$ be slots at integer nodes $s_i\ge0$. Then
\[
  \bigl|\det(\text{row of slot }(s_i,\tau_i))_{i<h}\bigr|\le h^h\Bigl(1+\frac2h\Bigr)^{h^2}\prod_{i<j}\operatorname{pf}(s_i,s_j),
\]
where $\operatorname{pf}(a,b)=|a-b|$ for $a\ne b$ and $\operatorname{pf}(a,a)=1$.
\end{lemma}

\begin{proof}
Let $\zeta$ be a primitive $h$-th root of unity and $r=1/h$. For $j<h$,
$\sum_{l<h}\zeta^{l(h-1)}(t+r\zeta^l)^j=h\,r\,jt^{j-1}$ (only the binomial term with $h\mid h-1+m$, i.e.\ $m=1$, survives;
this needs $h\ge2$). Since $hr=1$, the derivative row at $t$ is the sum
$\sum_l\zeta^{l(h-1)}(\text{value row at }t+r\zeta^l)$. By multilinearity the determinant is a sum of at most
$h^{\#\mathrm{der}}$ Vandermonde determinants at points $z_i$ with $|z_i-s_i|\le r$, each with a coefficient of modulus
$1$. Since the nodes are integers, each Vandermonde factor satisfies
$|z_j-z_i|\le|s_j-s_i|+2r\le(1+2r)\operatorname{pf}(s_i,s_j)$.
So the determinant is at most $h^h(1+2/h)^{h(h-1)/2}\prod_{i<j}\operatorname{pf}(s_i,s_j)$.
\end{proof}

Apply \eqref{eq:andreief} to $H^{(T)}$, indexing the rank-one terms by triples (node, type). A term with two equal
row vectors vanishes, so only assignments $f$ whose \emph{x-slots} are distinct contribute; in particular each node
carries at most two of the $h$ indices. With Lemma~\ref{lem:slotdet} for both determinants,
\begin{equation}\label{eq:pattern}
  |\det H^{(T)}|\le\frac{K_h^2}{h!}\sum_{f}\ \prod_{i<h}b(s_i)\prod_{i<j}\operatorname{pf}(s_i,s_j)^2,\qquad K_h:=h^h(1+2/h)^{h^2},
\end{equation}
where $f$ runs over the assignments $i\mapsto(s_i,\tau_i)$ with distinct x-slots and $s_i$ is the node of $f(i)$.

\subsection{A lattice energy inequality}

\begin{lemma}[P\'olya-type inequality]\label{lem:polya}
For $M\ge1$ and real $\nu_0,\dots,\nu_{M-1}$,
\[
  \sum_{i\ne j}\nu_i\nu_j\log|i-j|\le\log2\sum_i\nu_i^2+\log M\Bigl(\sum_i\nu_i\Bigr)^2 .
\]
\end{lemma}

\begin{proof}
For $M=1$ the claim reads $0\le\log2\,\nu_0^2$. Let $M\ge2$. We follow P\'olya's criterion for positive definite
functions \cite{Polya1949} in a finite form. Let $f(0)=\log M+\log2$, $f(d)=\log(M/d)$ for $1\le d\le M$, and $f(d)=0$
for $d>M$. Then $f$ is convex on
$\ZZ_{\ge0}$ (this uses $\log(k-1)+\log(k+1)\le2\log k$ for $2\le k<M$, $f(M-1)-2f(M)+f(M+1)=\log\frac M{M-1}\ge0$,
and $f(0)-2f(1)+f(2)=0$ when $M\ge2$).
Hence $f(d)=\sum_{k=1}^{M}\Delta^2f(k-1)\,(k-d)_+$ with $\Delta^2f\ge0$. Each triangle kernel $(k-|i-j|)_+$ is the
number of windows of length $k$ containing both $i$ and $j$, hence positive semidefinite. So
$\sum_{i,j}\nu_i\nu_jf(|i-j|)\ge0$, which rearranges to the claim.
\end{proof}

For a reference sequence $m=(m_s)_{s\ge0}$ with finite support, $\sum_sm_s=h$, put
\[
  P_m(s):=\sum_{s'\ne s}m_{s'}\log|s-s'|,\qquad E_m:=\sum_{s\ne s'}m_sm_{s'}\log|s-s'|.
\]

\begin{lemma}[discrete weak duality]\label{lem:dual}
Let $m$ be as above with $\sum_sm_s^2\le2h$, and let $\lambda'\in\RR$. Put $g(s):=\log b(s)+2P_m(s)-\lambda'$. For every assignment $f$ in \eqref{eq:pattern},
\[
  \prod_ib(s_i)\prod_{i<j}\operatorname{pf}(s_i,s_j)^2\le\exp\Bigl(-E_m+\lambda'h+8h\log2+\sum_ig(s_i)\Bigr).
\]
\end{lemma}

\begin{proof}
Let $n_s:=\#\{i:s_i=s\}\le2$, so $\sum_sn_s=h$ and $\sum_sn_s^2\le2h$. Pairs at the same node contribute
$\operatorname{pf}=1$, so $2\sum_{i<j}\log\operatorname{pf}(s_i,s_j)=\sum_{s\ne s'}n_sn_{s'}\log|s-s'|$. With $\nu:=n-m$ (total mass $0$),
\[
  \sum_{s\ne s'}n_sn_{s'}\log|s-s'|=-E_m+2\sum_sn_sP_m(s)+\sum_{s\ne s'}\nu_s\nu_{s'}\log|s-s'| ,
\]
and by Lemma~\ref{lem:polya} the last sum is at most $\log2\sum\nu_s^2\le2\log2(\sum n_s^2+\sum m_s^2)\le8h\log2$.
Finally $\sum_sn_sP_m(s)=\sum_iP_m(s_i)$.
\end{proof}

\subsection{Discretization}
Given a step density $\mu$ supported in $[0,N]$, let $m_s:=h\int_{(s-1)/h}^{s/h}\mu$ for $s\ge1$ and $m_0:=0$ (cell
averages). Then $\sum m_s=h$, $0\le m_s\le2$ (so $\sum m_s^2\le2h$), and $m_s=0$ for $s>Nh$.

\begin{lemma}\label{lem:discrete}
With these $m_s$, for all $h\ge2$ and $s\ge1$:
\begin{enumerate}
\item $\log b(s)\le\log h+\log s-h\log h+h\,V(s/h)$;
\item $P_m(s)\le h\,U^\mu(s/h)+h\log h+4+2\log(Nh+1)$;
\item $E_m\ge h^2E(\mu)+h^2\log h-4h\bigl(1+\log(Nh+1)\bigr)$.
\end{enumerate}
\end{lemma}

\begin{proof}
(i) By concavity of $\log$, $\log(c+\frac12)\ge\int_c^{c+1}\log y\,dy=\psi(c+1)-\psi(c)$ for $c\ge0$, so
$\log D(s)=\sum_{k<h}\log(s+k+\frac12)\ge\psi(s+h)-\psi(s)=h\bigl(\psi(x+1)-\psi(x)\bigr)+h\log h$ with $x=s/h$. Insert this into
$\log b(s)=s\log W+\log h+\log s-\log D(s)$.

(ii) We have $h\,U^\mu(s/h)=\sum_{s'}\int_{I_{s'}}\log|s-u|\,\rho(u)\,du-h\log h$, where $I_{s'}=[s'-1,s']$ and
$\rho(u):=\mu(u/h)$; this density has mass $m_{s'}$ on $I_{s'}$, total mass $h$, and is bounded by $2$. For
$s'<s$ and $u\in I_{s'}$, $\log|s-u|\ge\log(s-s')$, so these cells contribute with the right sign. For $s'=s$ the
discrete term is $0$ and $\int_{I_s}(-\log|s-u|)\,\rho\le2$. For $s'=s+d$, $d\ge1$, the deficit is at most
$2\int_{d-1}^{d}(\log d-\log u)\,du=2\bigl(1-(d-1)\log\tfrac d{d-1}\bigr)\le2/d$. Summing over $d\le Nh$ gives (ii).

(iii) Similarly $h^2E(\mu)+h^2\log h=\sum_{s,s'}\iint_{I_s\times I_{s'}}\log|u-u'|\,\rho(u)\rho(u')$. The diagonal cells
contribute $\le0$ because $|u-u'|\le1$ there. For $|s-s'|=d\ge1$ we have $\log|u-u'|\le\log(d+1)\le\log d+1/d$. Since $m_{s'}\le2$, summing
$m_sm_{s'}/d$ over the pairs gives at most $\sum_sm_s\cdot2\cdot2\sum_{d\le Nh}1/d\le4h(1+\log(Nh+1))$.
\end{proof}

\subsection{Proof of Theorem~\ref{thm:size}}
Fix a certificate and $h\ge2$. Take $m$ as above and $\lambda':=h\lambda+h\log h$. By Lemma~\ref{lem:discrete}, with
$\epsilon_h:=4+2\log(Nh+1)$ and $x=s/h$,
\begin{gather}
  g(s)\le h\bigl(F(x)-\lambda\bigr)+\log h+\log s+2\epsilon_h,\label{eq:gbound}\\
  -E_m+\lambda'h\le h^2\bigl(\lambda-E(\mu)\bigr)+O(h\log h).\notag
\end{gather}
\emph{Near and far nodes.} Let $S_0:=Xh$. Since each node carries at most two indices,
$\sum_ig(s_i)\le2\sum_{1\le s\le S_0}g(s)_++\sum_{i:\,s_i>S_0}g(s_i)$. Summing over $f$ in \eqref{eq:pattern} with
Lemma~\ref{lem:dual} therefore gives
\[
  |\det H^{(T)}|\le\frac{K_h^2}{h!}\,e^{-E_m+\lambda'h+8h\log2+2\sum_{s\le S_0}g(s)_+}\Bigl(3S_0+3\sum_{s>S_0}e^{g(s)}\Bigr)^h .
\]
\emph{Near nodes.} Each $x=s/h\le X$ lies in a cell $[k/L,(k+1)/L]$ with $k<XL$, where $F\le S_k$; each cell contains at
most $h/L+1$ points $s/h$. By \eqref{eq:gbound},
$\sum_{s\le S_0}g(s)_+\le h(h/L+1)\sum_k(S_k-\lambda)_++S_0(\log h+\log S_0+2\epsilon_h)$.

\emph{Far nodes.} Put $\delta:=-\max\bigl(1/X+\log W,\ (\log X+X\log W-\lambda)/X\bigr)>0$. For $x\ge N$,
$U^\mu(x)\le\log x$, and $V(x)\le x\log W-\log x$ (because $1-x\log(1+1/x)\le1/(x+1)\le\log(1+1/x)$). Hence, using
$\log x\le\log X+(x-X)/X$, for $x\ge X$
\begin{align*}
  F(x)-\lambda&\le\log x+x\log W-\lambda\\
  &\le(\log X+X\log W-\lambda)+(x-X)(1/X+\log W)\le-\delta x .
\end{align*}
So $e^{g(s)}\le h\,s\,e^{2\epsilon_h}e^{-\delta s}\le\frac2\delta h\,e^{2\epsilon_h}e^{-\delta s/2}$ for $s>S_0$, and the
far sum is $e^{O(\log h)}$.

\emph{Assembly.} $\log(K_h^2/h!)=O(h\log h)$, and $\log\bigl(3S_0+3\sum_{s>S_0}e^{g(s)}\bigr)=O(\log h)$. Collecting,
\[
  \log|\det H^{(T)}|\le h^2\Bigl(-E(\mu)+\lambda+\frac2L\sum_k(S_k-\lambda)_+\Bigr)+O(h\log h)
\]
uniformly in $T$, and the theorem follows from $|\Delta_h(\theta)|\le\sup_T|\det H^{(T)}|$. \qed

\begin{remark}
Let $C(W):=\sup\{E(\rho)+\int V\rho\}$ over densities $\rho$ on $[0,\infty)$ with compact support, $0\le\rho\le2$ and
$\int\rho=1$ (the constrained equilibrium constant). For such $\rho$, $E(\rho-\mu)\le0$ \cite[Lemma~I.1.8]{SaffTotik},
so $E(\rho)+\int V\rho\le-E(\mu)+\int F\rho\le-E(\mu)+\lambda+2\int_0^X(F-\lambda)_+\le\Phi$, using $\rho\le2$, $F<\lambda$
beyond $X$ and $F\le S_k$ on the cells. Hence
$\Phi\ge C(W)$. In \cite{Chi2Paper} the finite-$h$ values $h^{-2}\log|\Delta_h(\chi_2(1/q))|$ for $q=3,5$, computed in
ball arithmetic up to $h=320$, were observed to approach $C(1/9)\approx-2.1003$ and $C(1/25)\approx-2.4565$, so for
$w>0$ Theorem~\ref{thm:size} seems essentially sharp for this construction. For $w<0$ it is not: it sees only $|w|$, and
numerically $|\Delta^{\Ti}_h(\Ti_2(1/q))|$ is smaller than $|\Delta^\chi_h(\chi_2(1/q))|$ by a factor of roughly
$e^{-h^2/q^2}$ for small $q$. Only the form with an explicit certificate is used.
\end{remark}

% ===========================================================================
\section{Certificates and the shift lemma}\label{sec:cert}
% ===========================================================================

\subsection{The shift lemma}
A dual certificate depends on the point only through $W$, and only through the term $x\log W$ of $V_W$. Lowering $W$
therefore tilts the external field, and the cap $\mu\le2$ turns the tilt into a gain.

\begin{lemma}[shift lemma]\label{lem:shift}
Let $0<W'\le W<1$ and $\tau:=\log(W/W')$. For every dual certificate for $W$ with exponent $\Phi$ and every $\delta>0$
there is a dual certificate for $W'$ with exponent at most $\Phi-\tau/4+\delta$.
\end{lemma}

\begin{proof}
Keep $\mu$, $X$, $N$, and put $\lambda':=\lambda-\tau/2$. Since $V_{W'}(x)=V_W(x)-\tau x$, the new $F'$ is $F-\tau x$.
Refine the grid by a factor $2r$: $L'=2rL=2m$ with $m=rL$, and for $k'<XL'$ put $S'_{k'}:=S_k-\tau k'/L'$ with
$k=\lfloor k'/(2r)\rfloor$; on $[k'/L',(k'+1)/L']\subseteq[k/L,(k+1)/L]$ the term $-\tau x$ is at most its value at the left
end, so condition (i) holds. For (ii), $1/X+\log W'\le1/X+\log W<0$, and
$\log X+X\log W'-\lambda'=(\log X+X\log W-\lambda)-(X-\frac12)\tau<0$. Finally
$(S_k-\lambda+\tau(\frac12-k'/L'))_+\le(S_k-\lambda)_++\frac\tau2(1-k'/m)_+$, each $k$ is repeated $2r$ times, and
$\sum_{k'\ge0}(1-k'/m)_+=(m+1)/2$, so
\[
  \Phi'\le-E(\mu)+\lambda-\frac\tau2+\frac2L\sum_k(S_k-\lambda)_++\frac{2}{L'}\cdot\frac\tau2\cdot\frac{m+1}2
  =\Phi-\frac\tau4+\frac{\tau}{4m}.
\]
Choose $r$ with $\tau/(4rL)\le\delta$.
\end{proof}

For $W=q_0^{-2}$ and $W'=q^{-2}$ the gain is $\tau/4=\frac12\log(q/q_0)$; the same holds for $W=4q_0^{-2}$, $W'=4q^{-2}$.
Without the cap ($\mu$ of arbitrary height) there would be no gain: the cap forces every admissible density to have
mean at least $\frac14$.

\subsection{The certificates used}

\begin{theorem}\label{thm:cert}
There are dual certificates with the following exponents:
\begin{center}
\begin{tabular}{@{}lll@{}}
\toprule
weight $W$ & $\Phi\le$ & used for\\
\midrule
$1/3^2$ & $-2.078$ & Thm.~\ref{thm:A}, odd $q$\\
$1/14^2$ & $-3.03$ & Thm.~\ref{thm:A}, even $q\ge14$\\
$1/6^2$, $1/8^2$ & $-2.56$, $-2.72$ & Thm.~\ref{thm:A}, $q=6,8$\\
$1/10^2$, $1/12^2$ & $-2.85$, $-2.95$ & Thm.~\ref{thm:A}, $q=10,12$\\
$1/6^2$ (finer) & $-2.5671$ & Thm.~\ref{thm:B}, all $q\ge6$\\
$1/5^2$ & $-2.44$ & Thm.~\ref{thm:B}, $q=5$\\
$4/175^2$ & $-4.003$ & Thm.~\ref{thm:C}, odd $q\ge175$\\
\bottomrule
\end{tabular}
\end{center}
\end{theorem}

The certificates are explicit rational data (Table~\ref{tab:cert}). The density $\mu$ is constant on the $J$ cells
$[js/L,(j+1)s/L)$ with dyadic heights $c_j\in[0,2]$ (denominator $2^{16}$) and exact mass $\sum_jc_js/L=1$. The bounds
$S_k$ are rationals. All checks are inequalities between explicit rational numbers and finitely many values of $\log$ at
rational points, and are verified as follows.
\begin{itemize}
\item \emph{Logarithms.} A table of dyadic intervals of width $2\cdot2^{-40}$ for $\log n$, $n\le n_{\max}$, each checked
  against the series of $\log(1-x)$ with an explicit tail bound, after reducing the argument by powers of $2$; $\log2$
  is enclosed to width $2^{-100}$. The weight enters through $\log W=-2\log q$, or $2\log2-2\log175$.
\item \emph{Potential.} By summation by parts,
  $2U^\mu(x)+V(x)=\sum_jD_j\,\psi(x-\alpha_j/L)+x\log W$, with $D_j=2(c_j-c_{j-1})$ at $\alpha_j=js$, plus the terms
  $\psi(x)$ and $-\psi(x+1)$ of $V$. On each grid cell $[k/L,(k+1)/L]$ every term is bounded above by an affine
  function: by the chord where the term is convex, and by a tangent (at the midpoint, or at the endpoint away from
  $\alpha_j/L$) where it is concave. Summing the affine bounds and taking the larger endpoint value bounds $F$ on the
  cell. All values are rational enclosures.
\item \emph{Energy.} $E(\mu)=\sum_{j,j'}c_jc_{j'}\iint\log|x-y|$ over pairs of cells, and the double integral over
  $[a_1,b_1]\times[a_2,b_2]$ equals $-\Psi(b_1-b_2)+\Psi(b_1-a_2)+\Psi(a_1-b_2)-\Psi(a_1-a_2)$ with
  $\Psi(u)=\frac12u^2\log|u|-\frac34u^2$. A rational lower bound $E_{\mathrm{lo}}\le E(\mu)$ is computed in interval
  arithmetic.
\item \emph{Tail.} $1/X+\log W<0$ and $\log X+X\log W<\lambda$, from the enclosures of the logarithms.
\end{itemize}
The certified exponent is the rational number $\Phi_{\mathbb Q}:=-E_{\mathrm{lo}}+\lambda+\frac2L\sum_k(S_k-\lambda)_+\ge\Phi$.

\begin{table}[ht]
\centering\small
\begin{tabular}{@{}lcccccl@{}}
\toprule
$W$ & $1/L$ & $[0,X]$ ($XL$ cells) & $J\times$ width & $N$ & log table & $\Phi_{\mathbb Q}$\\
\midrule
$1/3^2$ & $1/100$ & $[0,3]$ ($300$) & $30\times\frac1{20}$ & $2$ & $804$ & $-2.0783260\ldots$\\
$1/14^2$ & $1/200$ & $[0,2]$ ($400$) & $40\times\frac1{50}$ & $1$ & $1204$ & $-3.0372394\ldots$\\
$1/6^2$ & $1/400$ & $[0,2]$ ($800$) & $54\times\frac1{50}$ & $2$ & $2404$ & $-2.5635463\ldots$\\
$1/8^2$ & $1/200$ & $[0,2]$ ($400$) & $50\times\frac1{50}$ & $1$ & $1204$ & $-2.7264892\ldots$\\
$1/10^2$ & $1/200$ & $[0,2]$ ($400$) & $46\times\frac1{50}$ & $1$ & $1204$ & $-2.8528812\ldots$\\
$1/12^2$ & $1/200$ & $[0,2]$ ($400$) & $46\times\frac1{50}$ & $1$ & $1204$ & $-2.9535971\ldots$\\
$1/6^2$ (finer) & $1/800$ & $[0,2]$ ($1600$) & $108\times\frac1{100}$ & $2$ & $4804$ & $-2.5671799\ldots$\\
$1/5^2$ & $1/200$ & $[0,2]$ ($400$) & $115\times\frac1{100}$ & $2$ & $1204$ & $-2.4435517\ldots$\\
$4/175^2$ & $1/1200$ & $[0,1]$ ($1200$) & $70\times\frac1{100}$ & $1$ & $4804$ & $-4.0033242\ldots$\\
\bottomrule
\end{tabular}
\caption{The dual certificates (decimal values truncated). The $q=3$ certificate is that of \cite{Chi2Paper}.}\label{tab:cert}
\end{table}

The data are generated by a script that optimises the density in floating point, rounds it to dyadic rationals, and
then replicates the interval arithmetic exactly in rational arithmetic. The script is not trusted: in the formalisation
every check above is a decidable statement about the data, proved by evaluation in the Lean kernel
(Section~\ref{sec:lean}). The certificates are not optimal. Finer independent certificates in ball arithmetic give,
for example, $\Phi\le-2.5704$ at $W=1/6^2$ and $\Phi\le-4.0005$ at $W=4/173^2$; they are used only in the remarks.

% ===========================================================================
\section{Nonvanishing: Lemma P}\label{sec:nonvan}
% ===========================================================================

\subsection{Statement}
In this section we work in case (A) with an integer $q\ge2$, so $z=1/q$, $w=q^{-2}$ and $\Delta_h=\Delta^\chi_h(X;q)$;
Sections~\ref{sec:B} and~\ref{sec:C} transfer the result to cases (B) and (C). For $r\in\{0,1,2\}$ put $c_r:=1/\prod_{a\in\{0,1,2\},a\ne r}(a-r)$, so $(c_0,c_1,c_2)=(\frac12,-1,\frac12)$. Let
$T(M)=\sum_{s\le M,\,s\text{ odd}}z^s/s^2$, with $T(M)=0$ for $M<1$, so that $S_k(q)=q^{2k+1}T(2k+1)$. Define three $3\times3$ rational matrices, indexed by $r,s\in\{0,1,2\}$:
\begin{align*}
  M_S(q)_{rs}&=\frac{c_r}{1-w}+\ind{r=s}\,4q^{2r+1}T(2r-1),\\
  M_L(q)_{rs}&=\frac{w\,c_r}{1-w}+\ind{r=s}\,4q^{2r+1}T(2r+1),\\
  M_D(q)_{rs}&=w\,c_r\Bigl(\frac{e_{rs}}{1-w}+\frac{2w}{(1-w)^2}\Bigr)-\ind{r=s}\,4\bigl(r+\tfrac12\bigr)q^{2r+1}T(2r+1),
\end{align*}
where $e_{rs}=\frac{11}2-2r$ if $r=s$ and $e_{rs}=(3-r-s)+\frac52$ if $r\ne s$.

\begin{theorem}[Lemma~P]\label{thm:lemmaP}
Let $q\ge2$ and let $\ell\ge7$ be a prime with $\ell\nmid q(q^2-1)$ such that $\det M_S(q)$, $\det M_L(q)$ and
$\det M_D(q)$ are $\ell$-adic units. Then $\Delta_{3\ell}(a/b)\ne0$ for all $a\in\ZZ$ and $b\ge1$ with $\ell\nmid b$.
More precisely, $|K_\ell\,\Delta_{3\ell}(a/b)|_\ell=1$, where $K_\ell\in\QQ^\times$ depends only on $\ell$ (not on $q$, $a$,
$b$).
\end{theorem}

\begin{proposition}[the block determinants]\label{prop:blocks}
With $Q=q^2$, $P_L(Q)=16200Q^3-10575Q^2-6352Q+905$ and $P_D(Q)=16200Q^3-51075Q^2+28523Q+420$,
\begin{gather*}
  \det M_S(q)=\frac{8Q^3(9Q+1)}{9(Q-1)},\qquad \det M_L(q)=\frac{8Q\,P_L(Q)}{2025(Q-1)},\\
  \det M_D(q)=-\frac{Q^2P_D(Q)}{135(Q-1)^2}
\end{gather*}
as rational functions of $q$. The numerators $9Q+1$, $P_L(Q)$ and $P_D(Q)$ do not vanish for real $Q\ge4$, nor for
real $Q\le-1$.
\end{proposition}

\begin{proof}
The identities follow by expanding the three $3\times3$ determinants, whose entries are explicit rational functions of
$q$ ($w=1/Q$, $T(1)=z$, $T(3)=z+z^3/9$, $T(5)=T(3)+z^5/25$, $T(-1)=0$); the formalisation carries out this expansion
exactly. For $Q\ge4$: $9Q+1>0$; $P_L(Q)\ge Q(16200Q^2-10575Q-6352)+905>0$ since $16200Q\ge10575+6352$; and
$P_D(Q)\ge Q^2(16200Q-51075)+28523Q+420>0$. For $Q=-x$ with $x\ge1$: $9Q+1<0$;
$P_L=-16200x^3-10575x^2+6352x+905<0$ because $10575x^2\ge10575x>6352x+905$; and
$P_D=-16200x^3-51075x^2-28523x+420<0$.
\end{proof}

\begin{corollary}[nonvanishing, case (A)]\label{cor:nonvanA}
For every integer $q\ge2$ and all $a\in\ZZ$, $b\ge1$, $\Delta^\chi_h(a/b;q)\ne0$ for infinitely many $h$.
\end{corollary}

\begin{proof}
By Proposition~\ref{prop:blocks} the three determinants are nonzero rationals. Every prime $\ell\ge7$ that divides none
of their numerators and denominators, nor $q(q^2-1)b$, satisfies the hypotheses of Theorem~\ref{thm:lemmaP}. There are
infinitely many such primes, and each gives $h=3\ell$.
\end{proof}

The rest of this section proves Theorem~\ref{thm:lemmaP}. Fix $q,\ell$ as there and $h=3\ell$; then $\ell^2\ge7\ell>2h$.
Let $\ZZ_{(\ell)}$ be the ring of $\ell$-integral rationals and $\bar{\cdot}\colon\ZZ_{(\ell)}\to\FF_\ell$ the reduction. Put
$\kappa:=(\ell-1)/2$. The indices $i<h$ fall into the classes $C_\gamma=\{\gamma,\gamma+\ell,\gamma+2\ell\}$, $\gamma<\ell$;
we write $i=\gamma+r\ell$ with $r\in\{0,1,2\}$. The class $\gamma$ has \emph{type} S if $\gamma<\kappa$, D if $\gamma=\kappa$, and L if
$\gamma>\kappa$. Note that $\ell\mid2i+1$ iff $i\in C_\kappa$; put $\delta_i:=\ind{i\in C_\kappa}$.

\subsection{Steps 1--2: the Lagrange Gram matrix}
For $r_0\in\QQ$ let $\Lambda_{r_0}:=A+r_0B$ and $G(r_0)_{ij}:=\Lambda_{r_0}(f_if_j)$, with the Lagrange basis $f_i=D/(u-x_i)$. By
\eqref{eq:gram}, $\det G(r_0)=(\det C_f)^2\Delta_h(r_0)$. Put
\[
  \mathcal E(r_0)_{ij}:=\nu_i\,G(r_0)_{ij}\,\ell^{-\delta_j},\qquad \nu_i:=\ell^2/D'(x_i).
\]
Then $\det\mathcal E(r_0)\ne0$ implies $\Delta_h(r_0)\ne0$; no value of $\det C_f$ is needed. As in the proof of
Lemma~\ref{lem:nine}, for $i\ne j$ we have $G(r_0)_{ij}=\mu(F_{ij})$ with $F_{ij}=u\prod_{m\ne i,j}(u-x_m)$, and
$G(r_0)_{ii}=\mu(Q_i)+4x_iD'(x_i)\bigl(S_i(q)-r_0q^{2i+1}\bigr)$. Hence
\[
  \mathcal E(r_0)=\mathcal E(0)+r_0\operatorname{diag}(d_i),\qquad d_i=-4\ell^{2-\delta_i}x_iq^{2i+1},\qquad v_\ell(d_i)=2 .
\]
If $\mathcal E(0)$ is $\ell$-integral and $\ell\nmid b$, then $\mathcal E(a/b)$ is $\ell$-integral with the same reduction as $\mathcal E(0)$.
So it suffices to show that $\mathcal E(0)$ is $\ell$-integral and $\det\overline{\mathcal E(0)}\ne0$.

\subsection{Step 3: splitting by progressions}
For $|v|<1$ and a polynomial $g$, $\sigma(v,g):=\sum_{j\ge0}v^jg(j)=\sum_b\Delta^bg(0)\,v^b/(1-v)^{b+1}$. Since $w$, $w^\ell$ and
$1-w^\ell\equiv1-w=(q^2-1)/q^2$ are $\ell$-units, we get the following. If $g$ has $\ell$-integral values on $\ZZ_{\ge0}$, then
$\sigma(w^\ell,g)$ is $\ell$-integral. If moreover $\bar g(j)=g_0(j)$ for all $j$, for a polynomial $g_0$ over $\FF_\ell$, then
$\overline{\sigma(w^\ell,g)}=\sum_b\Delta^bg_0(0)\,\bar w^b/(1-\bar w)^{b+1}$ (use $\bar w^\ell=\bar w$). Splitting $t\ge1$ into the
residue classes modulo $\ell$,
\[
  \mu(F)=\sum_{\gamma<\ell}w^{t_\gamma}\,\sigma\bigl(w^\ell,\ j\mapsto F'(t_\gamma+j\ell)\bigr),
\]
where the \emph{progression of $\gamma$} is $\{t\ge1:\ell\mid2t+2\gamma+1\}$ and $t_\gamma$ is its least element: $t_\gamma=\kappa-\gamma$
(type S), $\kappa-\gamma+\ell$ (type L), and $\ell$ (type D). On the progression of $\gamma$, $\ell\mid t-x_m$ iff $m\in C_\gamma$,
and $\ell\mid t$ iff $\gamma=\kappa$. Since $\ell^2>2h$, all these divisibilities are simple.

\subsection{Step 4: integrality and block triangularity}
We have $D'(x_i)=\prod_{m\ne i}(m-i)=\ell^2\,c_r^{-1}\,U_i$ with $U_i:=\prod_{m\notin C_\gamma}(x_i-x_m)$ an $\ell$-unit, so
$\nu_i=c_rU_i^{-1}$ is a unit. Let $\gamma_i,\gamma_j$ be the classes of $i,j$ and consider the contribution of the
progression of $\gamma$.

\emph{Off-diagonal entries.} At $t$ on the progression of $\gamma$, the factors of $F_{ij}$ divisible by $\ell$ number
$N_\gamma=3-\ind{\gamma=\gamma_i}-\ind{\gamma=\gamma_j}+\ind{\gamma=\kappa}$, and every term of $F'_{ij}(t)$ keeps at least $N_\gamma-1$ of
them. So the contribution has valuation $\ge N_\gamma-1-\delta_j\ge0$, with equality only if $\gamma=\gamma_i=\gamma_j$ (inside a
diagonal block) or $\gamma=\gamma_i\ne\gamma_j=\kappa$ (a non-D row against a D column).

\emph{Diagonal entries.} Off its own progression, $|t-x_i|_\ell=1$ and
$Q_i'(t)=\bigl(g_i'(t)(t-x_i)-g_i(t)+g_i(x_i)\bigr)/(t-x_i)^2$ with $g_i$ having at least three roots $\ell$-adically close to $t$,
so $v_\ell(Q_i'(t))\ge2$ and the contribution has valuation $\ge2-\delta_i\ge1$. The pole term
$4\ell^{2-\delta_i}x_iS_i(q)$ is integral, since $v_\ell(x_i)=\delta_i$ and the denominators $s^2$ of $S_i$ have $v_\ell(s)\le1$.
The own-progression contributions are integral by Step~5.

Hence $\mathcal E(0)$ is $\ell$-integral. Ordering the classes with $\kappa$ last, $\overline{\mathcal E(0)}$ is block upper triangular,
with $3\times3$ diagonal blocks $\overline{\mathcal E(0)}_{C_\gamma}$, and each diagonal block receives only its own progression
and the pole terms.

\subsection{Step 5: the diagonal blocks}
Fix a class $\gamma$ and $i=\gamma+r\ell$, $j=\gamma+s\ell$. On the progression of $\gamma$ we have $t\equiv x_i\equiv x_\gamma\pmod\ell$.

For $r\ne s$ let $c:=3-r-s$ be the rank of the third class member, and for $r=s$ let $c_1,c_2$ be the ranks of the
other two members; we write $x_{(c)}:=x_{\gamma+c\ell}$. For a polynomial $P$ let $Q[P]:=(P(u)-P(x_i))/(u-x_i)$, so that
$Q[PR]=Q[P]\,R+P(x_i)\,Q[R]$ and $Q_i=Q[g_i]$.

\emph{Types S and L.} If $r\ne s$, then
$F_{ij}=(u-x_{(c)})\cdot u\prod_{m\notin C_\gamma}(u-x_m)$, and modulo $\ell$, $F'_{ij}(t)\equiv t\prod_{m\notin C_\gamma}(t-x_m)\equiv x_iU_i$.
If $r=s$, write $g_i=P\cdot R$ with $P=(u-x_{(c_1)})(u-x_{(c_2)})$ and $R=u\prod_{m\notin C_\gamma}(u-x_m)$. Then
$Q_i'=Q[P]'R+Q[P]R'+P(x_i)Q[R]'$ with $Q[P]=u+x_i-x_{(c_1)}-x_{(c_2)}$. At $t$ we have $Q[P](t)\equiv0\pmod\ell$, and
$\ell^2\mid P(x_i)$, so $Q_i'(t)\equiv R(t)\equiv x_iU_i\pmod\ell$. In both cases $\nu_i(\cdot)\equiv c_r\bar x_\gamma$, a constant along
the progression. Summing with Step~3, and using $\bar w^\kappa=\bar q^{-(\ell-1)}=1$ and $\bar w^\ell=\bar w$, the progression
contributes $\bar w^{t_\gamma}c_r\bar x_\gamma/(1-\bar w)$, where $\bar w^{t_\gamma}=\bar w^{-\gamma}$ (S) or $\bar w^{1-\gamma}$ (L). For
the pole term only the odd $s=\ell s'$ survive in $\ell^2S_i$, and $q^\ell\equiv q$, so
$\overline{\ell^2S_i}=\bar q^{2i+1}\,T(M)$ with $M$ the largest odd number $\le(2i+1)/\ell$: $M=2r-1$ (S) or $2r+1$ (L). Since
$\bar q^{2i+1}=\bar q\,\bar w^{-\gamma}\bar w^{-r}$, the diagonal block is
\[
  \overline{\mathcal E(0)}_{C_\gamma}=\bar x_\gamma\,\bar w^{-\gamma}\,\overline{M_{\mathrm{type}(\gamma)}(q)} .
\]
\emph{Type D} ($\gamma=\kappa$). Here $x_i=-\ell(r+\frac12)$, the progression is $t=\ell(1+j')$, $j'\ge0$, and we put
$y=t/\ell+\frac12=j'+\frac32$; then $t=\ell(y-\frac12)$ and $t-x_{(c)}=\ell(y+c)$. Let
$R:=\prod_{m\notin C_\kappa}(u-x_m)$, so $R(t)\equiv R(x_i)=U_i\pmod\ell$. If $r\ne s$, two roots of $F_{ij}=u(u-x_{(c)})R$ are close
to $t$, and $F'_{ij}(t)=(2t-x_{(c)})R(t)+t(t-x_{(c)})R'(t)\equiv\ell(2y+c-\frac12)\,U_i\pmod{\ell^2}$. If $r=s$, write
$g_i=PR$ with $P=u(u-x_{(c_1)})(u-x_{(c_2)})$. Then $\ell^3\mid P(x_i)$; $Q[P](t)$ is a quadratic form in
$t,x_i,x_{(c_1)},x_{(c_2)}$, all divisible by $\ell$, so $Q[P](t)\equiv0\pmod{\ell^2}$; and $Q[P]'(t)=2t+x_i-x_{(c_1)}-x_{(c_2)}=\ell(2y+\frac52-2r)$ (using $c_1+c_2=3-r$).
Hence $Q_i'(t)\equiv\ell(2y+\frac52-2r)\,U_i\pmod{\ell^2}$. With $\nu_i=c_rU_i^{-1}$ and the column factor $\ell^{-1}$, the
residues are $c_r(2y+c-\frac12)=c_r(2j'+e_{rs})$ for $r\ne s$ and $c_r(2y+\frac52-2r)=c_r(2j'+e_{rr})$ for $r=s$. Summing
$c_r(2j'+e_{rs})$ against $\bar w^{1+j'}$ over $j'\ge0$ gives $\bar w\,c_r\bigl(e_{rs}/(1-\bar w)+2\bar w/(1-\bar w)^2\bigr)$.
The pole term $4\ell x_iS_i=-4(r+\frac12)\ell^2S_i$ reduces to
$-4(r+\frac12)\bar q^{2r+1}T(2r+1)$. So $\overline{\mathcal E(0)}_{C_\kappa}=\overline{M_D(q)}$.

\subsection{Step 6: conclusion}
By Steps 4 and 5,
\[
  \det\overline{\mathcal E(0)}=\prod_{\gamma\ne\kappa}\bigl(\bar x_\gamma\bar w^{-\gamma}\bigr)^3\,\overline{\det M_{\mathrm{type}(\gamma)}(q)}\cdot\overline{\det M_D(q)} .
\]
For $\gamma\ne\kappa$, $\bar x_\gamma=-(2\gamma+1)/2\ne0$, and the three determinants are units by hypothesis. So
$\det\overline{\mathcal E(a/b)}=\det\overline{\mathcal E(0)}\ne0$ in $\FF_\ell$, hence $\det\mathcal E(a/b)\ne0$ and
$\Delta_{3\ell}(a/b)\ne0$. Moreover $\mathcal E(a/b)$ is $\ell$-integral with a unit determinant, and
$\det\mathcal E(a/b)=\prod_i\nu_i\cdot\ell^{-3}\cdot(\det C_f)^2\,\Delta_{3\ell}(a/b)$, since exactly three columns
($j\in C_\kappa$) carry the factor $\ell^{-1}$. The factor $K_\ell:=\ell^{-3}(\det C_f)^2\prod_i\nu_i$ depends only on $\ell$.
This proves Theorem~\ref{thm:lemmaP}. \qed

\begin{remark}
The same analysis works for every $h$ with $h\bmod\ell\le\kappa$ and three or more atoms per class. One then needs a
fourth block type with four atoms, and an explicit prime-gap theorem to find suitable $\ell$ near $h/3$ for all large
$h$ (we have not written this out). For irrationality, infinitely many $h$ suffice, and $h=3\ell$ avoids both complications. With at
most two atoms per class the analogous block constants of type S vanish identically in our computations, so primes
$\ell>h/2$ did not give a usable nonvanishing statement.
\end{remark}

% ===========================================================================
\section{\texorpdfstring{Theorem~\ref{thm:A}: $\chi_2(1/q)$}{Theorem A: chi2(1/q)}}\label{sec:A}
% ===========================================================================

In this section we are in case (A): $w=q^{-2}$, $\theta=\chi_2(1/q)$, $\Delta_h=\Delta^\chi_h(X;q)$ with an integer $q\ge3$.
It remains to bound the prime $2$. We write $L:=\lfloor\log_2h\rfloor$ and use $|\cdot|_2$ coefficientwise on $\QQ[X]$.

\subsection{The prime 2 at odd \texorpdfstring{$q$}{q}}

\begin{lemma}[Lemma~10]\label{lem:twoA}
Let $q\ge3$ be odd and $\nu:=v_2(q^2-1)\ge3$. For every $h\ge1$ and $k\ge0$,
\[
  \bigl|[X^k]\Delta_h\bigr|_2\le2^{(\nu-2)\beta_h+h(L+6)} .
\]
\end{lemma}

For $q=3,5$ ($\nu=3$) this is \cite[Lemma~5.3]{Chi2Paper}; the proof there uses $\nu=3$ only in the bound for the entries
of type (L) below. We work with the anchored basis, for which $\det(\Lambda_X(a_ib_j))=\Delta_h(X)$.

\begin{lemma}\label{lem:fdseq}
Let $T\colon\ZZ_{\ge0}\to\QQ$ satisfy $T(l+1)=Q\,T(l)+e\,(s+2l)^{-2}$ with $v_2(e)\ge0$, $v_2(T(0))\ge0$ and $s$ odd, and
put $\eta:=Q-1$, $g(l):=(s+2l)^{-2}$. Then $\Delta^mT(0)=\eta^mT(0)+e\sum_{k<m}\eta^{m-1-k}\Delta^kg(0)$ and
$v_2(\Delta^kg(0))\ge k+v_2(k!)$. Consequently:
\begin{enumerate}
\item if $v_2(\eta)\ge2$, then $v_2(\Delta^mT(0))\ge v_2(m!)+m-2-\lfloor\log_2m\rfloor$ for $m\ge1$;
\item if $v_2(\eta)\ge1$, then $v_2(\Delta^mT(0))\ge m-1$ for $m\ge1$.
\end{enumerate}
\end{lemma}

\begin{proof}
$\Delta T=\eta T+eg$, and the formula follows by iterating. By induction on $k$,
$\Delta^kg(0)=(-2)^kk!\,H_k(s)/P_k(s)$ with $P_k(s)=\prod_{a\le k}(s+2a)$ odd and $H_k(s)=\sum_{a\le k}(s+2a)^{-1}$
$2$-integral, so $v_2(\Delta^kg(0))\ge k+v_2(k!)$. (i) The term $k=0$ has valuation $\ge2(m-1)$, which is enough. For
$k\ge1$, using $k-1-\lfloor\log_2k\rfloor\le v_2(k!)$ and $v_2(m!)\le m-1$, each term has valuation at least $2(m-1-k)+2k-1-\lfloor\log_2k\rfloor\ge v_2(m!)+m-2-\lfloor\log_2m\rfloor$, and
$v_2(\eta^mT(0))\ge2m$ is larger. (ii) Each term has valuation at least $(m-1-k)+k$, and $v_2(\eta^mT(0))\ge m$.
\end{proof}

\begin{proof}[Proof of Lemma~\ref{lem:twoA}]
We show that each entry satisfies
\begin{equation}\label{eq:twoentry}
  v_2\bigl(\Lambda_X(a_ib_j)\bigr)\ge h-(i+j+1)-(\nu-2)(i+j+1-h)_+-L-6 .
\end{equation}
\emph{Entries of type (L), $\pi:=i+j+1-h\ge1$.} Then $a_ib_j=D\cdot P$ with $P=\prod_{h-j\le m<i}(u-x_m)$, so $a_ib_j\varphi=uP$
is a polynomial: $B(a_ib_j)=0$ and $A(a_ib_j)=\mu(uP)$. We have $uP=2^{-\pi}H(2u)$ with $H\in\ZZ[T]$ of degree $\pi$.
Since $\Delta^b[(2t)^n](0)=2^nb!\,S(n,b)$ (Stirling numbers of the second kind) vanishes for $n<b$,
\begin{equation}\label{eq:Lnewton}
  v_2\bigl(\Delta^b[(uP)'](0)\bigr)\ge1-\pi+b+v_2(b!).
\end{equation}
By \eqref{eq:Omega}, $v_2(\Omega_b)=-\nu(b+1)$. By \eqref{eq:newton} and $v_2(b!)\ge b-1-\lfloor\log_2b\rfloor$,
\[
  v_2(A(a_ib_j))\ge\min_{b<\pi}\bigl(1-\pi-\nu-(\nu-1)b+v_2(b!)\bigr)\ge-(\nu-1)\pi-2-\lfloor\log_2\pi\rfloor ,
\]
which implies \eqref{eq:twoentry}.

\emph{Entries of type (U), $N:=h-1-i-j\ge0$.} Then $a_ib_j\varphi=u/\prod_{l=0}^{N}(u-\xi_l)$ with $\xi_l:=x_{i+l}=-c/2-l$,
$c:=2i+1$. Since $\xi_l-\xi_m=m-l$, its partial fractions are a constant (which $\mu$ kills) plus
$\sum_l(-1)^l\binom Nl\xi_l/(N!\,(u-\xi_l))$. By Lemma~\ref{lem:pole}, both $A(a_ib_j)$ and $B(a_ib_j)$ have the form
\begin{equation}\label{eq:Uentry}
  \pm\frac4{N!}\sum_l(-1)^l\binom Nl\xi_lT(l)=\pm\frac4{N!}\Delta^N[\xi_lT(l)](0),\qquad T(l)=s_{i+l}\ \text{or}\ c_{i+l}.
\end{equation}
Both satisfy the recursion of Lemma~\ref{lem:fdseq} with $Q=q^2$: $S_{i+l+1}=q^2S_{i+l}+(2i+3+2l)^{-2}$ ($e=1$,
$s=2i+3$) and $q^{2(i+l+1)+1}=q^2\cdot q^{2(i+l)+1}$ ($e=0$); here $v_2(Q-1)=\nu\ge2$. From
$\Delta^N[l\,T(l)](0)=N\,\Delta^{N-1}[T(\cdot+1)](0)$, and since $T(\cdot+1)$ satisfies a recursion of the same kind,
Lemma~\ref{lem:fdseq}(i) gives
\begin{multline*}
  v_2\Bigl(\frac4{N!}\Delta^N[\xi_lT(l)](0)\Bigr)\ge2-v_2(N!)\\
  +\min\Bigl(-1+v_2(N!)+N-2-\lfloor\log_2N\rfloor,\ v_2(N!)+N-3-\lfloor\log_2(N-1)\rfloor\Bigr)
\end{multline*}
for $N\ge2$, which is at least $N-1-\lfloor\log_2N\rfloor$. For $N\le1$ the entry is $\pm4\,\xi_0T(0)$ or
$\pm4\bigl(\xi_0\Delta T(0)-T(1)\bigr)$, of valuation $\ge1$, since $v_2(\xi_0)=-1$ and $T(0),T(1),\Delta T(0)$ are
$2$-integral. This implies \eqref{eq:twoentry}.

\emph{Assembly.} Since $(i+j+1-h)_+\le\rho_i+\kappa_j$, \eqref{eq:twoentry} gives
$|\Lambda_X(a_ib_j)|_2\le\bigl(2^{i+1+(\nu-2)\rho_i+L+6-h}\bigr)\bigl(2^{j+(\nu-2)\kappa_j}\bigr)$, and
Lemma~\ref{lem:leibniz} over $\QQ[X]$ gives
$|[X^k]\Delta_h|_2\le2^{\sum_i(i+1)+\sum_jj+(\nu-2)\beta_h+h(L+6)-h^2}=2^{(\nu-2)\beta_h+h(L+6)}$.
\end{proof}

So for odd $q$ the prime $2$ costs $\frac{\nu-2}4\log2$ per $h^2$, and with Theorem~\ref{thm:oddprimes} the arithmetic
constant is
\begin{equation}\label{eq:alphaodd}
\begin{split}
  \alpha(q)&=\frac{133}{60}+\frac14\log\odd(q^2-1)-\frac13\log q+\frac{\nu-2}4\log2\\
  &=\frac{133}{60}-\frac12\log2+\frac14\log(q^2-1)-\frac13\log q .
\end{split}
\end{equation}

\subsection{The prime 2 at even \texorpdfstring{$q$}{q}: a Gauss bound}
For even $q$ the lemma above has nothing to work with: $1-w$ is a $2$-adic unit. The $2$-adic size of $q$ itself enters
instead, through the $q$-adic gain of Lemma~\ref{lem:qadic}.

\begin{lemma}[Lemma~$10_0$]\label{lem:ten0}
Let $q\ge4$ be even and $e:=v_2(q)$. For every $h\ge1$ and $k\ge0$,
$\bigl|[X^k]\Delta_h\bigr|_2\le2^{\lfloor h^2/2\rfloor-e\,a_h}$.
\end{lemma}

\begin{proof}
First, $\|\Delta^{\cR}_h\|_2\le2^{\lfloor h^2/2\rfloor}$. Indeed, in the anchored basis over $\cR[X]$ (the Gram identity
\eqref{eq:gramR} with $\det C=1$) every entry has $\|\cdot\|_2\le2^{|i+j+1-h|}$. For type (L), $B^{\cR}=0$ and
$A^{\cR}=\mu^{\cR}(uP)$; the coefficient of $u^d$ in $uP=2^{-\pi}H(2u)$ has $v_2\ge d-\pi$, and $\mu^{\cR}(u^d)$ is an integer
combination of the $\Omega_b$, which have Gauss norm $1$. For type (U), the residue weights
$\xi_l(-1)^l/(l!\,(N-l)!)$ have $v_2\ge-1-v_2(N!)\ge-N$, against the factor $4$ and the pole factors
$q^{2k+1}X-S_k(q)$ of Gauss norm $\le1$ ($S_k$ has odd denominators). Since
$|i+j+1-h|\le|i-\frac{h-1}2|+|j-\frac{h-1}2|$, Lemma~\ref{lem:leibniz} gives $\|\Delta^{\cR}_h\|_2\le2^{2\sum_i|i-(h-1)/2|}$,
and $2\sum_{i<h}|i-\frac{h-1}2|=\lfloor h^2/2\rfloor$.

Let $c=[X^k]\Delta_h=\ev_q(f)$. Since $q^2-1$ is odd, \eqref{eq:clear} gives $|c|_2=|\ev(\cl_{\beta_h}f)|_2$. Every monomial of
$\cl_{\beta_h}f$ has degree $\ge a_h$ (Lemma~\ref{lem:qadic}), and $|q|_2=2^{-e}$, so
$|c|_2\le\|\cl_{\beta_h}f\|_2\,2^{-ea_h}\le\|f\|_2\,2^{-ea_h}$.
\end{proof}

So Lemma~\ref{lem:ten0} costs $(\frac12-\frac e3)\log2$ per $h^2$; it includes the $2$-adic share of the $q$-adic gain.

\subsection{The prime 2 at even \texorpdfstring{$q$}{q}: the weighted Gauss valuation}
Lemma~\ref{lem:ten0} loses twice: the Leibniz expansion in the anchored basis does not see the cancellation between its
terms, and the Gauss norm forgets that $q$ is $2$-adically small. Both losses disappear if we bound the $2$-adic size of the
\emph{reflected} slot expansion of Lemma~\ref{lem:qadic} directly, giving $q$ the weight $e$.

For $e\ge1$ define $v_e$ on $\QQ[[q]][X]$ by $v_e\bigl(\sum c_{n,j}q^nX^j\bigr):=\inf_{n,j}\bigl(v_2(c_{n,j})+en\bigr)$. It is
ultrametric and satisfies $v_e(FG)\ge v_e(F)+v_e(G)$.

\begin{lemma}[Lemma~$10_e$]\label{lem:tene}
Let $q\ge4$ be even and $e:=v_2(q)$. Put $m_e:=e-\frac{2(e+1)^2}{3e+8}$. Then for every $h\ge2$ and $k\ge0$,
\[
  -\log_2\bigl|[X^k]\Delta_h\bigr|_2\ge v_e\bigl(\psi_q(\Delta^{\cR}_h)\bigr)\ge m_eh^2-O(h\log h),
\]
with an implied constant depending only on $e$. So the prime $2$ costs $-m_e\log2=(c(e)-\frac e3)\log2$ per $h^2$, where
$c(e)=\frac{6-4e}{3(3e+8)}$: $c(1)=\frac2{33}$, $c(2)=-\frac1{21}$, $c(3)=-\frac2{17}$.
\end{lemma}

\begin{proof}
\emph{Step 1: reduction.} Let $f=[X^k]\Delta^{\cR}_h$ and $c=\ev_q(f)$. As $q^2-1$ is odd,
$v_2(c)=v_2(\ev(\cl_{\beta_h}f))$, and for a polynomial $N=\sum_nc_nq^n$ we have $v_2(N(q))\ge\min_n(v_2(c_n)+en)=v_e(N)$.
In $\QQ[[q]]$, $\cl_{\beta_h}f=(q^2-1)^{\beta_h}\psi_q(f)$ and $v_e(q^2-1)=0$, so $v_2(c)\ge v_e(\psi_q(f))\ge v_e(\psi_q(\Delta^{\cR}_h))$.

\emph{Step 2: the slot expansion.} By the proof of Lemma~\ref{lem:qadic}, for every $M$, the coefficients of $q^nX^j$ with
$n<2M$ in $\psi_q(\Delta^{\cR}_h)$ and in $\det H_M$ agree, where $H_M$ is the sum of the rank-one matrices
\begin{gather*}
  -q^{2m}\varphi'(-m)\,v_mv_m^{\mathsf T},\quad -q^{2m}\varphi(-m)\,v'_mv_m^{\mathsf T},\quad
  -q^{2m}\varphi(-m)\,v_mv_m'^{\mathsf T}\qquad(m<M),\\
  -4q^{2k+1}\omega_k(X+\chi_2(q))\,p_kp_k^{\mathsf T}\qquad(k<h),
\end{gather*}
with value rows $v_m=((-m)^i)_i$, derivative rows $v'_m=(i(-m)^{i-1})_i$ and pole rows $p_k=(x_k^i)_i$. So it suffices to
bound $v_e(\det H_M)$ from below uniformly in $M$, and by \eqref{eq:andreief} and $v_2(h!)\le h$ it suffices to bound each
Andreief term, $\prod_ic_{f(i)}\det X_f\det Y_f$, with pairwise distinct row keys in $X_f$ and in $Y_f$.

\emph{Step 3: the factors.} $D(-m)=\prod_k(k+\frac12-m)=2^{-h}\cdot(\text{odd})$, and $D'(-m)$ is a sum of products of $h-1$
such factors. Hence $v_2(\varphi(-m))\ge h$ and $v_2(\varphi'(-m))\ge h$. Next $v_2(4\omega_k)=1-v_2(k!)-v_2((h-1-k)!)\ge-h$,
and $v_e(X+\chi_2(q))\ge0$ because $\chi_2(q)=\sum q^{2j+1}/(2j+1)^2$ has $2$-integral coefficients. So a node slot at $m$
contributes $v_e\ge h+2em$ and a pole slot at $k$ contributes $v_e\ge-h+e(2k+1)$.

\emph{Step 4: the minors.} Each Andreief term attaches to every $i<h$ one point $z_i$, a node $-m$ or a pole $x_k$, and
$X_f$, $Y_f$ are both determinants with rows $(z_i^j)_j$ (value rows) or $(jz_i^{j-1})_j$ (derivative rows) at these same
points. Distinct keys allow at most two rows at a node (one value, one derivative) and one at a pole, and each rank-one
term has at most one derivative row among its two, so $X_f$ and $Y_f$ together have at most $I$ derivative rows, where
$I$ is the number of node slots. We use the following bound. \emph{Let $R$ have value rows at the points $z_i$, $i\notin U$,
and derivative rows at $z_i$, $i\in U$, and suppose $|z_i-z_{i'}|_2\ge2^{-\Lambda}$ whenever $i\in U$ and $z_{i'}\ne z_i$. Then
$|\det R|_2\le\prod_{i<i',\,z_i\ne z_{i'}}|z_i-z_{i'}|_2\cdot2^{\Lambda|U|}$.} Indeed $\det R=(\partial_UV)(z)$ with
$V=\prod_{i<i'}(X_{i'}-X_i)\in\ZZ[X_0,\dots,X_{h-1}]$ and $\partial_U=\prod_{i\in U}\partial/\partial X_i$, since only row $i$
involves $X_i$. By the product rule $\partial_UV$ is an integer combination of the products of all linear factors
$X_{i'}-X_i$ except those in $c(U)$, for injective maps $c$ from $U$ to the pairs with $i\in c(i)$. At $z$, such a product
vanishes unless every pair with equal points lies in $c(U)$, and then it is at most the product over pairs with distinct
points times $2^{\Lambda|U|}$ (one removed factor per $i\in U$, each of size $\ge2^{-\Lambda}$). The ultrametric inequality
gives the claim. We apply it with $\Lambda:=\lfloor\log_2M_f\rfloor$, where $M_f$ is the largest node of the term
(if there is no node slot, or only node $0$, there is no derivative row with another node to compare with, and we take
$\Lambda=0$, $M_f=1$): node
differences are nonzero integers of absolute value $\le M_f$, and node--pole differences $k+\frac12-m$ have $|\cdot|_2=2$.

\emph{Step 5: counting.} Let the term have $I$ node slots and $P=h-I$ pole slots. Then:
\begin{itemize}
\item prices: each node carries at most two slots, so $\sum2m\ge2\sum_{j<I}\lfloor j/2\rfloor+(2M_f-h)_+\ge\frac{I(I-1)}2-I+(2M_f-h)_+$;
  the poles are distinct, so $\sum(2k+1)\ge P^2$;
\item differences: for $L'\ge1$ let $\mathrm{cnt}(a)=\#\{l<L':2^{l+1}\mid a\}\le v_2(a)$. Over ordered pairs of node slots
  at distinct nodes, $\sum\mathrm{cnt}(m-m')\ge\sum_{l<L'}\bigl(I^2/2^{l+1}-2I\bigr)=I^2(1-2^{-L'})-2L'I$, by
  Cauchy--Schwarz over the residues modulo $2^{l+1}$ (at most $2I$ ordered pairs share a node). Likewise the pole pairs
  give at least $P^2(1-2^{-L'})-L'P$, and the $2IP$ ordered node--pole pairs give exactly $-2IP$;
\item derivatives: at most $I\le h$ rows lose $\Lambda$ each, and
  $h\lfloor\log_2M\rfloor\le e(2M-h)_++h(2\lfloor\log_2h\rfloor+1)$ for all $M\ge1$ (if $M\le h^2$ use
  $\lfloor\log_2M\rfloor\le2\lfloor\log_2h\rfloor+1$; if $M>h^2$ then $h\log_2M\le2M-h$).
\end{itemize}
Adding Steps 3--5 with $L'=\lceil\log_2h\rceil$, the term has
\[
  v_e\ge hI-hP+e\Bigl(\frac{I^2}2+P^2\Bigr)+I^2+P^2-2IP-O(h\log h)=h^2f_e(I/h)-O(h\log h),
\]
where $f_e(x)=e(\frac{x^2}2+(1-x)^2)+(2x-1)+(2x-1)^2=(\frac{3e}2+4)x^2-(2e+2)x+e$. Its minimum on $[0,1]$ is at
$x^*=\frac{2e+2}{3e+8}$ and equals $m_e$. This bound is uniform in $M$, which proves the lemma; and
$-m_e=\frac{2(e+1)^2}{3e+8}-e=c(e)-\frac e3$.
\end{proof}

\begin{remark}
The minimising slot sets are intervals: nodes $0,1,\dots$ and poles $0,\dots,P-1$, with about a third of the slots at
nodes when $e=1$. A similar count with $e=0$ (the plain Gauss valuation, applied to the polynomial
$\cl_{\beta_h}f$, whose $q$-degree is $O(h^2)$, so that only nodes $M=O(h^2)$ and losses $\log M=O(\log h)$ occur)
gives the cost $\frac14\log2$ in place of the $\frac12\log2$ of Lemma~\ref{lem:ten0}; we do not use this. Exact computations up to $h=256$ suggest that $c(e)$ is the true
asymptotic constant, approached slowly from above; we do not need this.
\end{remark}

\subsection{Proof of Theorem~\ref{thm:A}}
Let $q\ge3$, $q\ne4$. Hypothesis (iii) of Theorem~\ref{thm:conditional} is Corollary~\ref{cor:nonvanA}. For (i) we use
Theorem~\ref{thm:oddprimes} with the $2$-adic bound of Lemma~\ref{lem:twoA}, \ref{lem:ten0} or~\ref{lem:tene}; for (ii),
Theorem~\ref{thm:size} with a certificate of Theorem~\ref{thm:cert}, shifted by Lemma~\ref{lem:shift} where needed, with
$\delta=10^{-3}$. Write $\log(q^2-1)<2\log q$ throughout.
\begin{itemize}
\item \emph{Odd $q\ge3$}, base $W=1/9$: by \eqref{eq:alphaodd},
  \begin{align*}
    \alpha+\beta&<\frac{133}{60}-\frac12\log2+\frac12\log q-\frac13\log q-2.078-\frac12\log\frac q3+\delta\\
    &\le\frac{133}{60}-2.078-\frac12\log2+\frac16\log3+\delta<-0.023 .
  \end{align*}
\item \emph{Even $q\ge14$}, Lemma~\ref{lem:ten0}, base $W=1/196$: here
  $\alpha=\frac{133}{60}+\frac12\log2+\frac14\log(q^2-1)-\frac13\log q$, so
  \[
    \alpha+\beta<\frac{133}{60}+\frac12\log2-3.03+\frac16\log14+\delta<\frac{133}{60}-3.03+\frac76\log2+\delta<-0.003.
  \]
\item \emph{$q=6,8,10,12$}, Lemma~\ref{lem:tene}, direct certificates: $\alpha_e(q)+\Phi$ with
  $\alpha_e(q)=\frac{133}{60}+\frac14\log(q^2-1)-\frac13\log q+c(e)\log2$ is at most $-0.0097$, $-0.2422$, $-0.2101$, $-0.3539$
  (with $e=1,3,1,2$ and $\Phi\le-2.56,-2.72,-2.85,-2.95$).
\end{itemize}
The numerical values follow from $0.6931471803<\log2<0.6931471808$ and elementary bounds such as $\log3<\frac{27}{17}\log2$
($3^{17}<2^{27}$); the formalisation proves each inequality this way. Theorem~\ref{thm:conditional} gives
Theorem~\ref{thm:A}. \qed

\begin{remark}
For even $q\ge14$ Lemma~\ref{lem:tene} would give larger margins; we use Lemma~\ref{lem:ten0} there because that is how
the formal proof is organised. The case $q=6$ is the thinnest: with $c(1)=\frac2{33}$ it needs $\Phi<-2.5502$, and the
plain bound $c=\frac14$ would not suffice. For $q=4$ see Section~\ref{sec:limits}.
\end{remark}

% ===========================================================================
\section{\texorpdfstring{Theorem~\ref{thm:B}: $\Ti_2(1/q)$}{Theorem B: Ti2(1/q)}}\label{sec:B}
% ===========================================================================

In this section we are in case (B): $w=-q^{-2}$, $\theta=\Ti_2(1/q)$, $\Delta_h=\Delta^{\Ti}_h(X;q)$ with an integer $q\ge3$.
By Proposition~\ref{prop:twist}, $\Delta^{\Ti}_h$ is the image of the twisted symbolic determinant, so
Theorem~\ref{thm:oddprimes} already bounds the odd primes, with $Y=q^2+1$. In terms of the formal identity
$\Delta^{\Ti}_h(X;q)=\Delta^\chi_h(-iX;iq)$, $\Ti_2(1/q)$ is $\chi_2$ at the non-real point $i/q$; Theorem~\ref{thm:A} says
nothing about it.

\subsection{The prime 2 at odd \texorpdfstring{$q$}{q}}

\begin{lemma}[Lemma~$10_1$]\label{lem:ten1}
Let $q\ge3$ be odd. For every $h\ge1$ and $k\ge0$, $\bigl|[X^k]\Delta^{\Ti}_h\bigr|_2\le2^{\beta_h}$, $\beta_h=\lfloor h^2/4\rfloor$.
\end{lemma}

\begin{proof}
We follow the proof of Lemma~\ref{lem:twoA}. Here $q^2+1\equiv2\pmod4$, so by \eqref{eq:Omega} with $y=iq$,
$v_2(\Omega_b)=-(b+1)$: the analogue of $\nu$ is $1$. For entries of type (L), \eqref{eq:Lnewton} gives
$v_2(A(a_ib_j))\ge\min_{b<\pi}(1-\pi+b+v_2(b!)-b-1)\ge-\pi$, and $B(a_ib_j)=0$. For entries of type (U), \eqref{eq:Uentry}
holds with $T(l)=S^{\Ti}_{i+l}(q)$ or $(-1)^{i+l}q^{2(i+l)+1}$, which satisfy the recursion of Lemma~\ref{lem:fdseq} with
$Q=-q^2$, $\eta=-(q^2+1)$, $v_2(\eta)=1$, and $2$-integral initial values. By Lemma~\ref{lem:fdseq}(ii), for $N\ge2$,
$v_2(\Delta^N[\xi_lT(l)](0))\ge\min(-1+N-1,\ N-2)=N-2$, so the entry has $v_2\ge2-v_2(N!)+N-2=s_2(N)\ge1$; for $N\le1$
the entry is $\pm4\xi_0T(0)$ or $\pm4(\xi_0\Delta T(0)-T(1))$, of valuation $\ge1$. So
$|\Lambda_X(a_ib_j)|_2\le2^{(i+j+1-h)_+}\le2^{\rho_i}2^{\kappa_j}$, and Lemma~\ref{lem:leibniz} gives $2^{\beta_h}$.
\end{proof}

So for odd $q$ the prime $2$ costs $\frac14\log2$, and the arithmetic constant is
\begin{equation}\label{eq:alphaTi}
  \alpha^{\Ti}(q)=\frac{133}{60}+\frac14\log\frac{q^2+1}2-\frac13\log q+\frac14\log2=\frac{133}{60}+\frac14\log(q^2+1)-\frac13\log q .
\end{equation}
Compared with \eqref{eq:alphaodd}, $\Ti_2(1/q)$ pays $\frac12\log2+\frac14\log\frac{q^2+1}{q^2-1}$ more than
$\chi_2(1/q)$ at odd $q$.

\subsection{The prime 2 at odd \texorpdfstring{$q$}{q}: the binomial basis}
With Lemma~\ref{lem:ten1} the margin at $q=5$ is positive ($+0.055$ with the certificate $\Phi\le-2.44$ of
Theorem~\ref{thm:cert}, $+0.043$ with the ball-arithmetic certificate of the remark after the proof of
Theorem~\ref{thm:B}), and no bound on the entries of the anchored basis can do better: their tropical maximum is
$\approx h^2/4$, and the true $2$-adic size of $\Delta^{\Ti}_h$ comes from cancellation. Exact computations give
$-\min_kv_2([X^k]\Delta^{\Ti}_h)\approx h^2/8$. The next lemma proves this leading constant by changing to the binomial
basis, where the cancellation becomes a factorisation.

\begin{lemma}[Lemma~M]\label{lem:M}
Let $q\ge3$ be odd. Then $n_2(h):=-\min_kv_2\bigl([X^k]\Delta^{\Ti}_h\bigr)$ satisfies
$n_2(h)\le\frac{h^2}8+O(h\log h)$. So the prime $2$ costs at most $\frac18\log2$ per $h^2$.
\end{lemma}

\emph{Notation.} $C_s(u)=u(u-1)\cdots(u-s+1)/s!$ is the binomial polynomial: it has degree $s$ and leading coefficient
$1/s!$, it is integer-valued on $\ZZ_{\ge0}$, and $C_{s+1}(u+1)=C_{s+1}(u)+C_s(u)$ (Pascal). For $F\colon\ZZ_{\ge0}\to\QQ$,
$\Delta^bF(0)$ is the iterated forward difference. Put $B_0:=16h+16$,
\[
  \mathbb M(F):=\sum_{b<B_0}\Omega_b\,\Delta^bF(0),\qquad f_k(t):=-\frac1{(t-x_k)^2},
\]
\[
  \tilde\theta:=\frac{4-\mathbb M(f_0)}{4q}\in\QQ .
\]
Here $\tilde\theta$ is an auxiliary rational number, not the target $\Ti_2(1/q)$. By \eqref{eq:newton},
$\mu(R)=\mathbb M(R')$ for every polynomial $R$ of degree $\le B_0$. (The written proof in our working notes uses the
$2$-adic limit $\theta_2$ of $\tilde\theta$ and the untruncated sum; the truncation keeps everything rational.)

\begin{proof}
\emph{Step 1: the binomial Gram matrix.} By \eqref{eq:gram} with $a_i=b_i=C_i$,
$\Delta^{\Ti}_h=(\prod_{i<h}i!)^2\det(\Lambda_X(C_iC_j))_{i,j<h}$. For $k<h$, Newton's formula gives
$C_i(x_k)=\sum_{n<h}\binom kn\mathcal A_{ni}$ with $\mathcal A_{ni}:=\Delta^n_k[C_i(x_k)](0)$. Since $B(g)=\sum_k\beta_kg(x_k)$
with $\beta_k=-4\omega_k(-1)^kq^{2k+1}$, we get $B(C_iC_j)=(\mathcal A^{\mathsf T}\mathcal O\mathcal A)_{ij}$ with
$\mathcal O_{nn'}:=\sum_{k<h}\beta_k\binom kn\binom k{n'}$. As $\Lambda_X=(A+\tilde\theta B)+(X-\tilde\theta)B$,
\[
  \Lambda_X(C_iC_j)=E_{ij}+(X-\tilde\theta)(\mathcal A^{\mathsf T}\mathcal O\mathcal A)_{ij},\qquad
  E_{ij}:=A(C_iC_j)+\tilde\theta B(C_iC_j).
\]
By the Schur complement, $\det(E+(X-\tilde\theta)\mathcal A^{\mathsf T}\mathcal O\mathcal A)=\det\mathcal N$ for the
$3h\times3h$ matrix
$\mathcal N=\begin{pmatrix}E&\mathcal A^{\mathsf T}&0\\0&1&-(X-\tilde\theta)\mathcal O\\-\mathcal A&0&1\end{pmatrix}$ over
$\QQ[X]$.

\emph{Step 2: the entries of $\mathcal A$: $|\mathcal A_{ni}|_2\le2^{2(i-n)}$, and $\mathcal A_{ni}=0$ for $n>i$.} Since
$x_{k+1}=x_k-1$, Pascal gives $C_i(x_{k+1})-C_i(x_k)=-C_{i-1}(x_{k+1})$, so $\mathcal A_{ni}=(-1)^nC_{i-n}(x_n)$ for $n\le i$
and $\mathcal A_{ni}=0$ for $n>i$. And $C_s(x_k)=\prod_{l<s}\bigl(-\frac{2k+2l+1}2\bigr)/s!$ has $|\cdot|_2\le2^{2s}$.

\emph{Step 3: the entries of $\mathcal O$: $|\mathcal O_{nn'}|_2\le2^{n+n'-1}$.} With $K=h-1$ and $Q=q^2$,
$D'(x_k)=(-1)^kk!(K-k)!$ gives $\beta_k=\frac{2q}{K!}\binom Kk(2k+1)Q^k$. From
$\sum_k\binom Kk\binom ksQ^k=\binom KsQ^s(1+Q)^{K-s}$ and $k\binom ks=s\binom ks+(s+1)\binom k{s+1}$, the number
$T_s:=\sum_k\beta_k\binom ks$ is $\frac{2q}{K!}$ times an integer combination of $\binom KsQ^s(1+Q)^{K-s}$ and
$2\binom K{s+1}Q^{s+1}(1+Q)^{K-s-1}$ (with coefficients $2s+1$ and $s+1$); as $v_2(1+Q)=1$ and $v_2(K!)\le K$,
$v_2(T_s)\ge1-s$. Finally $\binom kn\binom k{n'}=\sum_{s\le n+n'}c_s\binom ks$ with $c_s\in\ZZ$ (the forward differences of an
integer-valued polynomial of degree $n+n'$), so $v_2(\mathcal O_{nn'})\ge1-n-n'$.

\emph{Step 4: the rational $\tilde\theta$.} As $f_{k+1}(t)=f_k(t+1)$, $\Delta^bf_{k+1}(0)=\Delta^bf_k(0)+\Delta^{b+1}f_k(0)$,
and $\Omega_{b-1}=(w^{-1}-1)\Omega_b$ ($b\ge2$), $\Omega_0=(w^{-1}-1)\Omega_1-1$ give
\[
  \mathbb M(f_{k+1})=-q^2\,\mathbb M(f_k)+\frac4{(2k+3)^2}+(-q^2-1)\,\Omega_{B_0}\,\Delta^{B_0}f_k(0).
\]
The pole values $D_k:=4S^{\Ti}_k(q)-4\tilde\theta(-1)^kq^{2k+1}$ satisfy $D_{k+1}=-q^2D_k+4/(2k+3)^2$, and
$D_0=\mathbb M(f_0)$ by the choice of $\tilde\theta$. Here $v_2(-q^2-1)=1$, $v_2(\Omega_{B_0})=-(B_0+1)$, and
$v_2(\Delta^{B_0}f_k(0))\ge B_0+2+v_2(B_0!)$ (the proof of Lemma~\ref{lem:fdseq}, as $f_k(t)=-4(2k+1+2t)^{-2}$), with
$v_2(B_0!)\ge B_0/2$. By induction on $k$, $|D_k-\mathbb M(f_k)|_2\le2^{-(8h+8)}$ for all $k$. Also
$v_2(\mathbb M(f_0))\ge1$, so $|\tilde\theta|_2\le2$.

\emph{Step 5: the entries of $E$: $|E_{ij}|_2\le2^{i+j+2-h}$.} Let $g=C_iC_j$ and $R_g=(gu)\operatorname{div}D$. By
definition $E_{ij}=\mu(R_g)+\sum_kg(x_k)\omega_kD_k$. Split $D_k=\mathbb M(f_k)+e_k$. The error terms have
$|g(x_k)\omega_ke_k|_2\le2^{2i+2j}\cdot2^h\cdot2^{-8h-8}\le2^{i+j+2-h}$ (Step~2, $v_2(\omega_k)\ge-h$, Step~4). The rest is
$\mu(R_g)+\sum_kg(x_k)\omega_k\mathbb M(f_k)=\mathbb M(F)$, where $F(t)=(gu/D)'(t)$ by partial fractions; here $(F_1/F_2)'(t)$
means $(F_1'(t)F_2(t)-F_1(t)F_2'(t))/F_2(t)^2$, so no analysis is needed. Since $|\Omega_b|_2=2^{b+1}$, it suffices to show
$|\Delta^bF(0)|_2\le2^{i+j+1-h-b}$ for $b<B_0$. Write $D(u)=2^{-h}P(2u)$ with $P(T)=\prod_{k<h}(T+2k+1)$, $P(0)$ odd. With
$U:=1-P/P(0)=T\,V$ and $Y:=P(0)^{-1}\sum_{m<M}U^m$ one has $PY=1-T^MV^M$, with $2$-integral coefficients, so
\[
  \frac{g(u)\,u}{D(u)}=2^h\,g(u)\,u\,Y(2u)+2^{h+M}\,\frac{g(u)\,u^{M+1}V(2u)^M}{P(2u)} .
\]
The polynomial part is $\sum_m2^{h+m}\sigma_mu^{m+1}g$ with $\sigma_m$ $2$-integral. For an integer-valued polynomial $G$ of
degree $\le d$, $|\Delta^bG'(0)|_2\le2^{d-b}$: write $G=\sum_sc_sC_s$ with $c_s\in\ZZ$, and use
$\Delta^bC_s'(0)=C_{s-b}'(0)=(-1)^{s-b-1}/(s-b)$ for $b<s$ (Pascal, differentiated), and $0$ for $b\ge s$. With
$d=m+1+i+j$ this gives $2^{i+j+1-h-b}$. For the tail take $M=B_0+2h+1$: at a node $t$, $P(2t)$ is odd; the numerator
$g(u)\,u^{M+1}V(2u)^M$ is an integer there, and its derivative has $|\cdot|_2\le2^{2h}$: the factor $u^{M+1}V(2u)^M$
has integer coefficients, $g(t)\in\ZZ$, and $|g'(t)|_2\le2^{i+j}\le2^{2h}$ by the derivative bound just proved (with
$b=0$, applied to $G(x)=g(t+x)$). So the tail contributes at most $2^{2h-h-M}\le2^{i+j+1-h-b}$.

\emph{Step 6: one Leibniz bound with potentials.} Put $\alpha=-\lfloor h/2\rfloor$, $\gamma=2-h+\lfloor h/2\rfloor$ and
$\beta=\alpha-\lfloor h/4\rfloor$, and give the rows of $\mathcal N$ the exponents $\max(i+\alpha,2i+\beta)$ ($E$-rows,
$i<h$), $\max(2n+\beta,n)$ (middle rows) and $-2n-\beta$ (bottom rows), and the columns the exponents
$\max(j+\gamma,2j+\beta)$ ($E$-columns, $j<h$), $-2n-\beta$ (middle columns) and $\max(2n+\beta,n)$ (right columns). By
Steps~2--5, and since the coefficients of $-(X-\tilde\theta)\mathcal O_{nn'}$ are $-\mathcal O_{nn'}$ and
$\tilde\theta\mathcal O_{nn'}$, every entry of $\mathcal N$ has coefficients of $2$-adic size at most
$2^{\text{row}+\text{column}}$. Lemma~\ref{lem:leibniz} bounds the coefficients of $\det\mathcal N$ by $2^{S_h}$, where $S_h$
is the sum of all the exponents. Writing $\max(x,y)=y+(x-y)_+$ and summing,
$2S_h\le4h(h-1)-4hr+p(p+1)+p'(p'+1)+2r(r+1)$ with $r=-\beta$, $p=\lfloor h/4\rfloor$, $p'=\gamma-\beta$, and this is at most
$\frac94h^2+13$. With $v_2(i!)\ge i-1-\log_2i$, $n_2(h)\le S_h-2\sum_{i<h}v_2(i!)\le\frac{h^2}8+O(h\log h)$.
\end{proof}

\begin{remark}
The crossover of the potentials at $h/4$ is the extremal configuration of the proof in our working notes: a quarter of
the slots at node levels, three quarters at pole levels. With the sharper bound $|E_{ij}|_2\le2^{i+j+1-h}$ of those notes,
the tropical value of the Leibniz bound for $\mathcal N$ equals their subset bound exactly (we checked this for
$h\le48$); the version above is weaker by $O(h)$. Exact computations at $q=3,5$ give $n_2(h)=h^2/8+O(h)$ for $h\le80$, so
the constant $\frac18$ appears to be sharp.
\end{remark}

\subsection{The prime 2 at even \texorpdfstring{$q$}{q}}

\begin{corollary}[Lemma~$10_e$ for $\Ti_2$]\label{cor:teneTi}
Let $q\ge4$ be even and $e=v_2(q)$. Then $-\log_2|[X^k]\Delta^{\Ti}_h|_2\ge m_eh^2-O(h\log h)$, with $m_e$ as in
Lemma~\ref{lem:tene}.
\end{corollary}

\begin{proof}
Let $c=[X^k]\Delta^{\Ti}_h=\ev^{\Ti}_q(f^{\Ti})$. Since $q^2+1$ is odd, $v_2(c)=v_2(\ev(\cl^-_{\beta_h}f^{\Ti}))\ge v_e(\cl^-_{\beta_h}f^{\Ti})$.
By the remark after \eqref{eq:clear}, $\cl^-_{\beta_h}f^{\Ti}$ and $\cl_{\beta_h}f$ have the same coefficients up to sign, so
$v_e(\cl^-_{\beta_h}f^{\Ti})=v_e(\cl_{\beta_h}f)\ge v_e(\psi_q(\Delta^{\cR}_h))$, which Lemma~\ref{lem:tene} bounds.
\end{proof}

\subsection{Nonvanishing}
Let $M^{\Ti}_S(q)$, $M^{\Ti}_L(q)$, $M^{\Ti}_D(q)$ be the matrices of Section~\ref{sec:nonvan} with $w=-q^{-2}$ and with $q^{2r+1}T(M)$
replaced by $(-1)^rq^{2r+1}T^\natural(M)$, where $T^\natural(M)=\sum_{s\le M,\,s\text{ odd}}\chi_4(s)q^{-s}/s^2$, so that
$(-1)^kq^{2k+1}T^\natural(2k+1)=S^{\Ti}_k(q)$. Their determinants are the right-hand sides of
Proposition~\ref{prop:blocks} at $Q=-q^2$; for instance $\det M^{\Ti}_S(q)=-8q^6(9q^2-1)/(9(q^2+1))$. By
Proposition~\ref{prop:blocks} (the case $Q\le-1$) they are nonzero rationals.

\begin{proposition}[nonvanishing, case (B)]\label{prop:nonvanB}
Theorem~\ref{thm:lemmaP} holds for $\Delta^{\Ti}_{3\ell}$ with $\ell\nmid q(q^2+1)$ and $M^{\Ti}_S$, $M^{\Ti}_L$, $M^{\Ti}_D$
in place of $M_S$, $M_L$, $M_D$. Hence for every integer $q\ge3$ and all $a\in\ZZ$, $b\ge1$, $\Delta^{\Ti}_h(a/b)\ne0$ for
infinitely many $h$.
\end{proposition}

\begin{proof}
The proof of Theorem~\ref{thm:lemmaP} uses $w$ only through the facts that $w$, $w^\ell$ and $1-w$ are $\ell$-adic units
and $\bar w^\ell=\bar w$, which hold for $w=-q^{-2}$ when $\ell\nmid q(q^2+1)$, and through two congruences in Step~5,
which change by signs. First, $\bar w^\kappa=(-1)^\kappa\bar q^{-(\ell-1)}=\chi_4(\ell)$ instead of $1$. Second, in
$\ell^2S^{\Ti}_i$ only the terms with $\ell\mid s$ survive, and $\chi_4(\ell s')=\chi_4(\ell)\chi_4(s')$. In the blocks of
types S and L both changes occur together, and the block becomes the unit $(-1)^{\kappa+\gamma}$ times
$\bar x_\gamma\bar q^{2\gamma}\,\overline{M^{\Ti}_{\mathrm{type}}(q)}$; in the block of type D the two signs cancel (in the pole
term $(-1)^i=(-1)^\kappa(-1)^r$), and the block is $\overline{M^{\Ti}_D(q)}$. So $\det\overline{\mathcal E(0)}$ is a unit
times the product of the three determinants, and Step~6 concludes as before. The last statement follows as in
Corollary~\ref{cor:nonvanA}.
\end{proof}

\subsection{Proof of Theorem~\ref{thm:B}}
Let $q\ge5$. Theorem~\ref{thm:size} applies with $W=|w|=q^{-2}$: the certificates see only $|w|$. For $q=5$ we use
Lemma~\ref{lem:M} and the certificate at $W=1/25$ of Theorem~\ref{thm:cert} ($\Phi\le-2.44$):
\[
  \frac{133}{60}+\frac14\log13-\frac13\log5+\frac18\log2-2.44=-0.0319\ldots .
\]
For $q\ge6$ we use the finer
certificate at $W=1/36$ of Theorem~\ref{thm:cert} ($\Phi\le-2.5671$) directly at $q=6$ and shifted by
Lemma~\ref{lem:shift} for $q\ge7$, $\beta=-2.5671-\frac12\log\frac q6+\delta$. With $\log(q^2+1)\le2\log q+q^{-2}$:
\begin{itemize}
\item $q=6$ ($e=1$, Corollary~\ref{cor:teneTi}): $\frac{133}{60}+\frac14\log37-\frac13\log6+\frac2{33}\log2-2.5671=-0.00294\ldots$;
\item odd $q\ge7$ (Lemma~\ref{lem:ten1}): by \eqref{eq:alphaTi},
  $\alpha+\beta\le\frac{133}{60}-2.5671+\frac12\log6+\frac1{4q^2}-\frac13\log q+\delta$, which decreases in $q$ and is
  $<-0.097$ at $q=7$;
\item even $q\ge8$ (Corollary~\ref{cor:teneTi}, $c(e)\le c(1)$): the same bound plus $\frac2{33}\log2$, $<-0.10$ at
  $q=8$.
\end{itemize}
With $\delta=10^{-4}$ all margins are negative, and Theorem~\ref{thm:conditional} gives Theorem~\ref{thm:B}. \qed

\begin{remark}
For odd $q\ge7$ Lemma~\ref{lem:ten1} suffices, and the formal proof uses it there; Lemma~\ref{lem:M} would give larger
margins. At $q=5$ a finer ball-arithmetic certificate ($\Phi\le-2.4516816$) gives $-0.0436$. The margin at $q=6$ is
about $-0.003$ per $h^2$ (only Theorem~\ref{thm:C} at $q=175$, about $-0.0014$, is thinner). It rests on the exact leading constant
$c(1)=\frac2{33}$ of Lemma~\ref{lem:tene}; with $c=\frac14$ the margin would be $+0.13$. A finer certificate in ball
arithmetic ($\Phi\le-2.5704$) gives $-0.0063$.
\end{remark}

% ===========================================================================
\section{\texorpdfstring{Theorem~\ref{thm:C}: $\chi_2(2/q)$}{Theorem C: chi2(2/q)}}\label{sec:C}
% ===========================================================================

In this section we are in case (C): $q\ge5$ odd, $\zeta=q/2$, $w=\zeta^{-2}=4q^{-2}$, $\theta=\chi_2(2/q)$ and
$\Delta_h=\Delta^\chi_h(X;\zeta)=\ev_\zeta(\Delta^{\cR}_h)$. This is the machine of Theorem~\ref{thm:A} read at the rational point
$\zeta$, and Theorem~\ref{thm:oddprimes} bounds the odd primes with $Y=q^2-4$. At $p=2$, however, $\zeta$ is not
$2$-integral: $v_2(\zeta)=-1$ and $v_2(\zeta^2-1)=-2$.

\subsection{The prime 2}

\begin{lemma}[Lemma~$10''$]\label{lem:tenZ}
For odd $q\ge3$, every $h\ge1$ and $k\ge0$, $\bigl|[X^k]\Delta_h\bigr|_2\le2^{n_2(h)}$ with
\[
  n_2(h)=4hm-3m^2+2m,\qquad m=\Bigl\lfloor\frac{2h+2}3\Bigr\rfloor,\qquad 3n_2(h)\le4h^2+4h+4 .
\]
So the prime $2$ costs at most $\frac43\log2$ per $h^2$.
\end{lemma}

\begin{proof}
Use the anchored basis. \emph{Type (L)}, $\pi=i+j+1-h\ge1$: $B=0$, and by \eqref{eq:Omega} with $y=\zeta$,
$\Omega_0=4/(q^2-4)$ and $\Omega_b=\zeta^2\Omega_0^{b+1}$ have $v_2(\Omega_b)=2b+[b=0]\cdot2\ge0$, so by \eqref{eq:Lnewton}
$v_2(A)\ge1-\pi\ge-\pi$. \emph{Type (U)}, $N=h-1-i-j\ge0$: by \eqref{eq:Uentry} and the triangle inequality, with
$v_2(\xi_l)=-1$, $v_2(\zeta^{2k+1})=-(2k+1)$ and $v_2(S_k(\zeta))\ge-2k$, each term $\xi_lT(l)$ has
$v_2\ge-2(i+l)-2\ge-2i-2N-2$, and with $v_2(N!)\le N$ the entry has $v_2\ge-2i-3N$. So $v_2(\Lambda_X(a_ib_j))\ge-c(i,j)$ with
$c(i,j)=i+j+1-h$ on (L) and $c(i,j)=2i+3(h-1-i-j)$ on (U).

Put $a=\lfloor(4h+2)/3\rfloor$ and
\[
  U_i=\begin{cases}-i,&3i<2h,\\ i-a,&3i\ge2h,\end{cases}\qquad
  V_j=\begin{cases}3h-3j,&3j<2h,\\ j+1-h+a,&3j\ge2h,\end{cases}
\]
integer versions of the potentials $u(x)=\max(-x,x-\frac43)$, $v(y)=\max(3-3y,y+\frac13)$. A check of the four cases
$3i\gtrless2h$, $3j\gtrless2h$ shows $c(i,j)\le U_i+V_j$ for all $i,j<h$. Lemma~\ref{lem:leibniz} with $\rho_i=2^{U_i}$,
$\kappa_j=2^{V_j}$ gives $|[X^k]\Delta_h|_2\le2^{\sum_i(U_i+V_i)}$, and the sum equals $4hm-3m^2+2m$. Its continuous version
is $h^2\int_0^1(u+v)=\frac43h^2$.
\end{proof}

The constant $\frac43$ cannot be lowered much by any $2$-adic argument for this machine:

\begin{proposition}[the price of the numerator]\label{prop:floor}
For odd $q$ and every $h\ge1$, $[X^h]\Delta_h=\pm4^h\zeta^{h^2}\prod_{k<h}x_k$, which has $v_2=h-h^2$. Hence every $d_h$
with $d_h\Delta_h\in\ZZ[X]$ has $v_2(d_h)\ge h^2-h$.
\end{proposition}

\begin{proof}
$[X^h]\Delta_h=\det(B(u^{i+j}))$, and $B(u^{i+j})=\sum_kb_kx_k^{i+j}$ with $b_k=-4\omega_k\zeta^{2k+1}$, so
$\det(B(u^{i+j}))=\prod_kb_k\cdot(\det V)^2$. Since $(\det V)^2=\prod_{k<l}(l-k)^2=\prod_k|D'(x_k)|$, this is
$\pm4^h\prod_kx_k\cdot\zeta^{\sum_k(2k+1)}=\pm4^h\zeta^{h^2}\prod_kx_k$, of valuation $2h-h^2-h$.
\end{proof}

For the numerator $3$ we observed numerically that the $3$-adic loss is exactly that of the leading coefficient, $\log3$
per $h^2$; at $p=2$ the measured loss (about $1.18\log2$ at $h=64$) is larger than $\log2$, presumably because of the
half-integer poles.

\subsection{Nonvanishing}

\begin{proposition}[nonvanishing, case (C)]\label{prop:nonvanC}
For every odd $q\ge5$ and all $a\in\ZZ$, $b\ge1$, $\Delta^\chi_h(a/b;q/2)\ne0$ for infinitely many $h$.
\end{proposition}

\begin{proof}
We transfer Theorem~\ref{thm:lemmaP} from an integer point by $\ell$-adic continuity. Let $\ell\ge7$ be a prime that divides
neither $30\,b\,q(q^2-4)$ nor the integers $9q^2+4$, $64P_L(q^2/4)$, $64P_D(q^2/4)$, which are nonzero by
Proposition~\ref{prop:blocks} ($q^2/4\ge4$). Fix $M\ge1$ and put $q':=q(\ell^M+1)/2$, an integer with $2q'\equiv q\pmod{\ell^M}$.
Then $4q'^2\equiv q^2$, so $\ell\nmid q'(q'^2-1)$, and by Proposition~\ref{prop:blocks} the numerators and denominators of
$\det M_\bullet(q')$ are congruent modulo $\ell$ to units times the integers above; so $q'$ and $\ell$ satisfy the
hypotheses of Theorem~\ref{thm:lemmaP}, and $|K_\ell\,\Delta^\chi_{3\ell}(r;q')|_\ell=1$ for $r=a/b$.

Now let $h=3\ell$ and let $f\in\cR$ be a coefficient of $\Delta^{\cR}_h$; by Lemma~\ref{lem:nine}, $\|f\|_\ell\le\ell^{2h\lev(\ell)}$.
The points $y_1=q/2$ and $y_2=q'$ are $\ell$-adic units with $y_1^2-1$, $y_2^2-1$ units and $y_1\equiv y_2\pmod{\ell^M}$,
so $q^{m_0}E^{m_1}$ takes values at $y_1$ and $y_2$ that are congruent modulo $\ell^M$, and
$|\ev_{y_1}(f)-\ev_{y_2}(f)|_\ell\le\ell^{2h\lev(\ell)-M}$. As $r$ is $\ell$-integral,
$|\Delta_h(r;q/2)-\Delta_h(r;q')|_\ell\le\ell^{2h\lev(\ell)-M}$. Taking $M$ so large that $|K_\ell|_\ell\,\ell^{2h\lev(\ell)-M}<1$,
we get $|K_\ell\,\Delta_h(r;q/2)|_\ell=1$, so $\Delta_h(r;q/2)\ne0$. Infinitely many primes $\ell$ qualify.
\end{proof}

\subsection{Proof of Theorem~\ref{thm:C}}
Let $q\ge175$ be odd. By Theorem~\ref{thm:oddprimes} and Lemma~\ref{lem:tenZ}, the arithmetic constant is
\[
  \alpha(q)=\frac{133}{60}+\frac14\log(q^2-4)-\frac13\log q+\frac43\log2 .
\]
Theorem~\ref{thm:size} with the certificate at $W=4/175^2$ ($\Phi\le-4.003$), shifted by Lemma~\ref{lem:shift} to $W=4/q^2$,
gives $\beta=-4.003-\frac12\log\frac q{175}+\delta$. With $\log(q^2-4)<2\log q$,
\[
  \alpha+\beta<\frac{133}{60}+\frac43\log2-4.003+\frac16\log175+\delta<-0.0012+\delta ,
\]
using $175^{104}<2^{775}$ and $\log2<0.6931471808$. With $\delta=10^{-4}$ and Proposition~\ref{prop:nonvanC},
Theorem~\ref{thm:conditional} gives Theorem~\ref{thm:C}. \qed

\begin{remark}[$q=173$]\label{rem:173}
Two independent ball-arithmetic computations (Arb), which are not included in this paper and not formalised, give
$\Phi\le-4.0005336$ and $\Phi\le-4.0004556$ at $W=4/173^2$. With either of them the margin at $q=173$ would be at most
$-0.00082$ (respectively $-0.00074$), and the argument would reach every odd $q\ge173$. We do not claim this: in the
certificate format of Section~\ref{sec:cert}, a grid of about $L=3200$ cells would be needed at $W=4/173^2$, which we
have not run. The value $q=171$ is out of reach for
Lemma~\ref{lem:tenZ} even with the exact equilibrium constant (numerically the margin is $+0.0016$).
\end{remark}

% ===========================================================================
\section{Limits of the method}\label{sec:limits}
% ===========================================================================

The values below are numerical unless stated otherwise. ``Margin'' means $\alpha+\Phi$ per $h^2$, with the proved
arithmetic constants and the best certificate (or the equilibrium constant $C(W)$, which no certificate can beat).
\begin{itemize}
\item \emph{$\chi_2(1/4)$.} With Lemma~\ref{lem:tene} ($c(2)=-\frac1{21}$) and the certificate $\Phi\le-2.3088$ the margin is
  $+0.090$, and the equilibrium constant does not change this. Closing it would need $c(2)<-0.177$, while exact
  computations give about $-0.025$ at $h=160$. So $q=4$ needs a different construction.
\item \emph{$\chi_2(1/2)$.} The margin is about $+0.52$ \cite[\S11.4]{Chi2Paper}; the analytic constant $C(\frac14)\approx-1.744$ is
  the obstruction.
\item \emph{$\Ti_2(1/q)$ for $q\le4$.} With Lemma~\ref{lem:M} the margin at $q=3$ is $+0.241$, and at $q=4$ (even, where
Lemma~\ref{lem:tene} already applies) it is $+0.121$. Improving the $2$-adic bound cannot close either gap: at $q=3$
removing the prime $2$ altogether ($\frac18\log2\approx0.087$) leaves $+0.154$, and at $q=4$ the even constant $c(2)$ would
have to drop far below what exact computations show, as for $\chi_2(1/4)$ above. The size bound sees only $|w|$, and the true determinants at $w<0$ are smaller
(Section~\ref{sec:analytic}), so a signed version of Theorem~\ref{thm:size} is the remaining lever; $q=2$ fails for $\chi_2$
already.
\item \emph{$\chi_2(2/q)$ below $173$.} Proposition~\ref{prop:floor} is proved: no $2$-adic lemma can give less than $\log2$
  per $h^2$ for this machine. With $c_2=\log2$ and the equilibrium constant, the margin is positive for odd $q\lesssim87$
  (it changes sign between $81$ and $101$); with the measured $2$-adic constant (about $1.18\log2$ at $h=64$) the sign
  change is between $121$ and $141$. So, numerically and assuming that $\frac{133}{60}$ is the true generic constant,
  this machine does not reach odd $q$ below about $87$, and the gap between $87$ and $175$ is the gap between the $2$-adic
  floor and Lemma~\ref{lem:tenZ} (plus the coarseness of our certificates).
\item \emph{Effectivity.} Every bound is effective except the prime number theorem in Theorem~\ref{thm:generic}, which
  is used without an error term. The theorems are not affected.
\end{itemize}

% ===========================================================================
\section{The formalisation}\label{sec:lean}
% ===========================================================================

\subsection{What is formalised}
The repository \url{https://github.com/telaviv/chi2-irrationality} \cite{Chi2Repo}, at commit \lean{cef469e} of branch
\lean{main} (full hash \lean{cef469ef846a2a3967d2d441e331254bd82898e2}), contains proofs of Theorems~\ref{thm:A},
\ref{thm:B} and~\ref{thm:C} in Lean~4 \cite{LeanFour} (version v4.35.0) with Mathlib \cite{Mathlib} (version v4.35.0
\cite{Mathlib4}). The same repository contains the formalisation of \cite{Chi2Paper} ($q=3,5$), on which the present
one builds, and an intermediate theorem ($\chi_2(1/q)\notin\QQ$ for $q\ge3$, $q\notin\{4,6,8,10,12\}$), which the proof of
Theorem~\ref{thm:A} reuses. It has 148 Lean files under \lean{Chi2Irrational\slash{}}, including the certificate data.
Paths below are relative to \lean{Chi2Irrational\slash{}}, and all declarations are in the namespace \lean{Chi2Irr}.
Every step of the proofs of Theorems~\ref{thm:A}--\ref{thm:C} is formalised, including the prime number theorem, the Andreief identity, the lattice
energy inequality, all dual certificates and the block determinants.

\subsection{The statements}
The definitions that the reader has to check are:
\begin{quote}\ttfamily\footnotesize
-- Defs.lean\\
noncomputable def chi2Term (z : $\mathbb{R}$) (n : $\mathbb{N}$) : $\mathbb{R}$ :=\\
\hspace*{2em}z \^{} (2 * n + 1) / (2 * (n : $\mathbb{R}$) + 1) \^{} 2\\
noncomputable def chi2 (z : $\mathbb{R}$) : $\mathbb{R}$ := $\sum$' n : $\mathbb{N}$, chi2Term z n\\[2pt]
-- TiQ/Defs.lean (namespace TiQ)\\
noncomputable def ti2Term (z : $\mathbb{R}$) (n : $\mathbb{N}$) : $\mathbb{R}$ :=\\
\hspace*{2em}(-1) \^{} n * z \^{} (2 * n + 1) / (2 * (n : $\mathbb{R}$) + 1) \^{} 2\\
noncomputable def ti2 (z : $\mathbb{R}$) : $\mathbb{R}$ := $\sum$' n : $\mathbb{N}$, ti2Term z n\\[2pt]
-- Q/Even/Defs.lean, Q/Even/Interface.lean (namespace Q)\\
def TheoremQ4 : Prop := $\forall$ q : $\mathbb{N}$, 3 $\le$ q $\to$ q $\ne$ 4 $\to$\\
\hspace*{2em}Irrational (chi2 (1 / (q : $\mathbb{R}$)))\\
theorem theoremQ4 : TheoremQ4\\[2pt]
-- TiQ/Interface5.lean (namespace TiQ)\\
def TheoremT5 : Prop := $\forall$ q : $\mathbb{N}$, 5 $\le$ q $\to$ Irrational (ti2 (1 / (q : $\mathbb{R}$)))\\
theorem theoremT5 : TheoremT5\\[2pt]
-- QZ/Interface.lean (namespace QZ)\\
theorem theoremQ4num : $\forall$ q : $\mathbb{N}$, Odd q $\to$ 175 $\le$ q $\to$\\
\hspace*{2em}Irrational (chi2 (2 / (q : $\mathbb{R}$)))
\end{quote}
So Theorem~\ref{thm:A} is \lean{Q.theoremQ4}, Theorem~\ref{thm:B} is \lean{TiQ.theoremT5} (which uses \lean{TiQ.theoremT}, the case $q\ge6$, from \lean{TiQ\slash{}Interface.lean}) and Theorem~\ref{thm:C} is
\lean{QZ.theoremQ4num}. Here $\sum'$ is Mathlib's unconditional sum (\lean{tsum}); \lean{summable\_chi2Term} and
\lean{TiQ.summable\_ti2Term} prove that the series converge absolutely for $|z|\le1$, so \lean{chi2 z} and \lean{ti2 z}
are the genuine values. \lean{Irrational x} is Mathlib's predicate $x\notin\QQ$, and $1/(q:\RR)$, $2/(q:\RR)$ are real
divisions after the cast $\NN\to\RR$. The statements do not involve the machine, so a wrong internal definition could
make a proof fail but could not change what is proved.

\subsection{Axioms}
The command \lean{lake env lean Axioms.lean} prints, for $72$ declarations (the three theorems above, the theorems of
\cite{Chi2Paper}, and the main ingredients of each),
\begin{quote}\ttfamily\small
'Chi2Irr.Q.theoremQ4' depends on axioms: [propext, Classical.choice, Quot.sound]\\
'Chi2Irr.TiQ.theoremT5' depends on axioms: [propext, Classical.choice, Quot.sound]\\
'Chi2Irr.QZ.theoremQ4num' depends on axioms: [propext, Classical.choice, Quot.sound]
\end{quote}
and the same list for each of the other $69$ (run on 2026-11-13 on the tree of the commit above; the output is
\lean{scripts\slash{}axioms\_t5.txt}). These are Lean's three standard axioms, which Mathlib also uses. The
development contains no \lean{sorry}, no added axiom, no \lean{native\_decide} and no \lean{implemented\_by}.

\subsection{Correspondence with the paper}
Tables~\ref{tab:leanshared}--\ref{tab:leanC} list the formal counterparts of the main steps.

\begin{table}[ht]
\centering\small
\begin{tabular}{@{}>{\raggedright\arraybackslash}p{0.27\textwidth}>{\raggedright\arraybackslash}p{0.69\textwidth}@{}}
\toprule
paper & Lean declaration (file) \\
\midrule
Lemma~\ref{lem:pole}, case (A) & \lean{pole\_series} (\lean{PoleSeries}); \lean{Lam\_eq}, \lean{aeval\_chi2\_Delta} (\lean{PoleExpansion})\\
Prop.~\ref{prop:criterion}, Thm.~\ref{thm:conditional} & \lean{irrational\_of\_poly\_approx} (\lean{Criterion}); \lean{irrational\_of\_bounds} (\lean{MainCond})\\
Section~\ref{sec:free} & \lean{map\_DeltaF} (\lean{Arith\slash{}Free}); \lean{det\_GramF} (\lean{Arith\slash{}Gauss})\\
Lemma~\ref{lem:nine} & \lean{lemma9\_DeltaF} (\lean{Arith\slash{}Lemma9})\\
Prop.~\ref{prop:classbound}, Lemma~\ref{lem:classcount} & \lean{classBound\_DeltaF} (\lean{Arith\slash{}ClassBound}); \lean{classBd\_le\_phi1} (\lean{Arith\slash{}ClassCount})\\
Thm.~\ref{thm:generic} & \lean{eventually\_primeSum\_phi1\_le}, \lean{eventually\_primeSum\_lev\_le} (\lean{Arith\slash{}PrimeSums}); \lean{PNT.theta\_div\_tendsto} (\lean{PNT\slash{}Theta})\\
Prop.~\ref{prop:twist} & \lean{TiQ.parity\_DeltaF}, \lean{TiQ.map\_DeltaFT}, \lean{TiQ.coeff\_clearET\_twF} (\lean{TiQ\slash{}Free})\\
Lemma~\ref{lem:poleorder}, \eqref{eq:clear} & \lean{EDeg\_DeltaF}, \lean{aeval\_clearE} (\lean{Arith\slash{}PoleOrder})\\
Lemma~\ref{lem:qadic} & \lean{ord\_DeltaP}, \lean{qdeg\_clearE\_DeltaF} (\lean{Arith\slash{}QAdic}); \lean{factorial\_smul\_det\_sum\_rankOne} (\lean{Analytic\slash{}Andreief})\\
Thm.~\ref{thm:oddprimes} & \lean{Q.arith\_q} (\lean{Q\slash{}Arith}); \lean{TiQ.arith\_t} (\lean{TiQ\slash{}Arith}); \lean{QZ.arith\_z} (\lean{QZ\slash{}Arith})\\
Def.~\ref{def:cert}, Thm.~\ref{thm:size} & \lean{DualCert}, \lean{DualCert.analytic} (\lean{Analytic\slash{}TheoremB}); \lean{TiQ.analytic\_t\_proof} (\lean{TiQ\slash{}Analytic}); \lean{QZ.DualCertW}, \lean{QZ.DualCertW.analyticZ} (\lean{QZ\slash{}TheoremB})\\
Lemma~\ref{lem:shift} & \lean{Q.DualCert.shift} (\lean{Q\slash{}Shift}); \lean{QZ.DualCertW.shift} (\lean{QZ\slash{}Shift})\\
Thm.~\ref{thm:cert} & \lean{Cert.CData.Valid}, \lean{Cert.CData.Phi\_le} (\lean{Cert\slash{}Instance}); \lean{QZ.CDataZ.Valid} (\lean{QZ\slash{}CertInstance})\\
Thm.~\ref{thm:lemmaP} & \lean{Delta\_eval\_ne\_zero\_of\_det} (\lean{NonVan\slash{}Gram}); \lean{lemmaP} (\lean{NonVan\slash{}LemmaP}); \lean{nonvan\_of} (\lean{NonVan\slash{}NonVanQ})\\
\bottomrule
\end{tabular}
\caption{Formal counterparts of the shared steps.}\label{tab:leanshared}
\end{table}

\begin{table}[ht]
\centering\small
\begin{tabular}{@{}>{\raggedright\arraybackslash}p{0.27\textwidth}>{\raggedright\arraybackslash}p{0.69\textwidth}@{}}
\toprule
paper & Lean declaration (file) \\
\midrule
Lemma~\ref{lem:twoA} & \lean{Q.twoAdic\_Delta}, \lean{Q.twoAdic\_odd\_proof} (\lean{Q\slash{}TwoAdic}); \lean{Q.V2\_entry} (\lean{Q\slash{}TwoAdicEntries})\\
Lemma~\ref{lem:ten0} & \lean{Q.GLeX\_DeltaF\_two}, \lean{Q.twoAdic\_even\_proof} (\lean{Q\slash{}TwoAdicEven})\\
Lemma~\ref{lem:tene} & \lean{Q.CoreDet.padicNorm\_detR\_le} (\lean{Q\slash{}Even\slash{}CoreDet}); \lean{Q.sum\_cntL\_ge}, \lean{Q.large\_node} (\lean{Q\slash{}Even\slash{}Count}); \lean{Q.twoAdic\_even\_sharp\_proof} (\lean{Q\slash{}Even\slash{}TwoAdicSharp})\\
Prop.~\ref{prop:blocks}, Cor.~\ref{cor:nonvanA} & \lean{Q.nonvan\_q\_proof} (\lean{Q\slash{}NonVan})\\
certificates & \lean{Cert.Q3.analytic} (\lean{Cert\slash{}Q3}); \lean{Cert.Q14.phi\_le} (\lean{Cert\slash{}Q14}); \lean{Q.cert\_newQ\_proof} (\lean{Q\slash{}Even\slash{}Certs}; data \lean{Cert\slash{}Q6}, \dots, \lean{Cert\slash{}Q12})\\
margins & \lean{Q.margin\_odd}, \lean{Q.margin\_even} (\lean{Q\slash{}Interface}, proofs in \lean{Q\slash{}Margin}); \lean{Q.margin\_newQ} (\lean{Q\slash{}Even\slash{}Interface}, \lean{Q\slash{}Even\slash{}Margin})\\
Theorem~\ref{thm:A} & \lean{Q.theoremQ} (\lean{Q\slash{}Interface}; odd $q$, even $q\ge14$); \lean{Q.theoremQ4} (\lean{Q\slash{}Even\slash{}Interface})\\
\bottomrule
\end{tabular}
\caption{Formal counterparts for Theorem~\ref{thm:A}.}\label{tab:leanA}
\end{table}

\begin{table}[ht]
\centering\small
\begin{tabular}{@{}>{\raggedright\arraybackslash}p{0.27\textwidth}>{\raggedright\arraybackslash}p{0.69\textwidth}@{}}
\toprule
paper & Lean declaration (file) \\
\midrule
case (B) of Lemma~\ref{lem:pole} & \lean{TiQ.pole\_seriesT} (\lean{TiQ\slash{}PoleSeries}); \lean{TiQ.Lam\_eqT}, \lean{TiQ.aeval\_ti2\_DeltaT} (\lean{TiQ\slash{}PoleExpansion})\\
Lemma~\ref{lem:ten1} & \lean{TiQ.V2\_fd\_seq1}, \lean{TiQ.V2\_entryT} (\lean{TiQ\slash{}TwoAdicEntries}); \lean{TiQ.twoAdic\_DeltaT}, \lean{TiQ.twoAdic\_odd\_t\_proof} (\lean{TiQ\slash{}TwoAdicNu1})\\
Cor.~\ref{cor:teneTi} & \lean{TiQ.twoAdic\_even\_t\_proof} (\lean{TiQ\slash{}TwoAdicEven})\\
Prop.~\ref{prop:nonvanB} & \lean{TiQ.block\_SL}, \lean{TiQ.block\_D} (\lean{TiQ\slash{}NonVan\slash{}BlockMat}); \lean{TiQ.lemmaP} (\lean{TiQ\slash{}NonVan\slash{}LemmaP}); \lean{TiQ.nonvan\_t\_proof} (\lean{TiQ\slash{}NonVan\slash{}NonVanT})\\
Thm.~\ref{thm:size} at $w<0$ & \lean{TiQ.abs\_WT}, \lean{TiQ.analytic\_t\_proof} (\lean{TiQ\slash{}Analytic})\\
certificate $W=1/36$ & \lean{TiQ.cert\_t6\_proof} (\lean{TiQ\slash{}Certs}; data and checks \lean{TiQ\slash{}Cert\slash{}})\\
Lemma~\ref{lem:M}, Steps 1--2 & \lean{TiQ.M.DeltaT\_eq\_Cb}, \lean{TiQ.M.BfunT\_Cb}, \lean{TiQ.M.HM\_eq}, \lean{TiQ.M.det\_block} (\lean{TiQ\slash{}M\slash{}Gram}); \lean{TiQ.M.Am\_zero}, \lean{TiQ.M.Am\_bound} (\lean{TiQ\slash{}M\slash{}Binom})\\
Lemma~\ref{lem:M}, Steps 3--5 & \lean{TiQ.M.Om\_bound} (\lean{TiQ\slash{}M\slash{}BoundB}); \lean{TiQ.M.fd\_deriv\_intval} (\lean{TiQ\slash{}M\slash{}Binom}); \lean{TiQ.M.MB\_poly}, \lean{TiQ.M.err\_bound}, \lean{TiQ.M.theta\_bound} (\lean{TiQ\slash{}M\slash{}Trunc}); \lean{TiQ.M.boundA} (\lean{TiQ\slash{}M\slash{}BoundA})\\
Lemma~\ref{lem:M}, Step 6 & \lean{TiQ.M.det\_PNz}, \lean{TiQ.M.Nblk\_entry}, \lean{TiQ.M.twoAdic\_DeltaT\_M}, \lean{TiQ.M.eventually\_n2M\_le} (\lean{TiQ\slash{}M\slash{}Final}); \lean{TiQ.twoAdic\_odd\_M} (\lean{TiQ\slash{}Interface5})\\
certificate $W=1/25$ & \lean{TiQ.cert\_t5} (\lean{TiQ\slash{}Interface5}; data and checks \lean{Cert\slash{}DataQ5L200}, \lean{Cert\slash{}Q5L200})\\
margins & \lean{TiQ.margin\_t5} (\lean{TiQ\slash{}Interface5}); \lean{TiQ.margin\_t6\_proof}, \lean{TiQ.margin\_t\_odd\_proof}, \lean{TiQ.margin\_t\_even\_proof} (\lean{TiQ\slash{}Margin})\\
Theorem~\ref{thm:B} & \lean{TiQ.irrational\_of\_boundsT} (\lean{TiQ\slash{}MainCond}); \lean{TiQ.theoremT} (\lean{TiQ\slash{}Interface}, $q\ge6$); \lean{TiQ.theoremT5} (\lean{TiQ\slash{}Interface5})\\
\bottomrule
\end{tabular}
\caption{Formal counterparts for Theorem~\ref{thm:B}.}\label{tab:leanB}
\end{table}

\begin{table}[ht]
\centering\small
\begin{tabular}{@{}>{\raggedright\arraybackslash}p{0.27\textwidth}>{\raggedright\arraybackslash}p{0.69\textwidth}@{}}
\toprule
paper & Lean declaration (file) \\
\midrule
case (C), \eqref{eq:spec} & \lean{QZ.DeltaZ}, \lean{QZ.map\_DeltaF\_Z} (\lean{QZ\slash{}Machine}); \lean{QZ.pole\_seriesZ}, \lean{QZ.aeval\_chi2\_DeltaZ} (\lean{QZ\slash{}PoleExpansion}); \lean{QZ.chi2\_two\_div} (\lean{QZ\slash{}Interface})\\
Lemma~\ref{lem:tenZ} & \lean{QZ.V2\_entryZ} (\lean{QZ\slash{}TwoAdicEntries}); \lean{QZ.twoAdic\_DeltaZ}, \lean{QZ.twoAdic\_num\_proof} (\lean{QZ\slash{}TwoAdic})\\
Prop.~\ref{prop:nonvanC} & \lean{QZ.lemmaP\_norm}, \lean{QZ.nonvan\_z\_proof} (\lean{QZ\slash{}NonVan})\\
certificate $W=4/175^2$ & \lean{QZ.Cert.Q175.valid} (\lean{QZ\slash{}CertQ175}; data \lean{QZ\slash{}CertDataQ175})\\
margin & \lean{QZ.margin\_z\_proof} (\lean{QZ\slash{}Margin})\\
Theorem~\ref{thm:C} & \lean{QZ.irrational\_of\_boundsZ} (\lean{QZ\slash{}MainCond}); \lean{QZ.theoremQ4num} (\lean{QZ\slash{}Interface})\\
\bottomrule
\end{tabular}
\caption{Formal counterparts for Theorem~\ref{thm:C}.}\label{tab:leanC}
\end{table}

\subsection{Differences from the written proof}\label{sec:lean-departures}
The formal proofs follow the written ones. The differences are these; none weakens a stated theorem.
\begin{itemize}
\item \emph{Organisation.} For odd $q$ and even $q\ge14$, \lean{Q.theoremQ4} calls the intermediate theorem
  \lean{Q.theoremQ}; for $q=6,8,10,12$ it uses Lemma~\ref{lem:tene}. The machines of cases (B) and (C) are separate copies
  of the machine of case (A) (\lean{TiQ\slash{}Machine}, \lean{QZ\slash{}Machine}), tied to the symbolic determinant by
  \lean{TiQ.map\_DeltaFT} and \lean{QZ.map\_DeltaF\_Z}; no file of the earlier formalisation was edited. Lemma~\ref{lem:pole},
  the size bound and Theorem~\ref{thm:lemmaP} are copied with the needed sign or weight changes.
\item \emph{Lower-order terms.} Lemma~\ref{lem:ten0} is formalised with $h^2/2+h$ in place of $\lfloor h^2/2\rfloor$;
  Lemma~\ref{lem:ten1} with weaker entry bounds ($v_2\ge-\pi-2-\lfloor\log_2\pi\rfloor$ on (L), $v_2\ge0$ on (U)), giving
  $n_2\le\beta_h+h(\lfloor\log_2h\rfloor+2)$; Lemma~\ref{lem:tene} with the explicit bound
  $n_2\le h-\lceil m_eh^2-h(\frac{3e}2+4\lfloor\log_2h\rfloor+4)\rceil$. The leading constants are those of the paper. The value
  $\beta_h=\lfloor h^2/4\rfloor$ is not formalised; the proofs use $4\beta_h\le h^2+2h$.
\item \emph{Shift and margins.} The shift lemma is formalised in the discretised form of Lemma~\ref{lem:shift}. The margins
  are proved with elementary bounds for logarithms, so the formal margins are slightly weaker: for instance $-0.00904$ at
  $q=6$ for Theorem~\ref{thm:A}, about $-0.0025$ at $q=6$ and $-0.019$ at $q=5$ for Theorem~\ref{thm:B}, and $-0.00126$ for Theorem~\ref{thm:C}.
\item \emph{Lemma~\ref{lem:M}.} The formal proof is Steps 1--6 above (rational $\tilde\theta$, binomial basis, the
  block matrix $\mathcal N$ and one Leibniz bound with potentials). These differ from our working notes, not from the
  Lean: the notes use the $2$-adic $\theta_2$ and a Cauchy--Binet expansion over $2h$ rows, with the sharper bound
  $|E_{ij}|_2\le2^{i+j+1-h}$; with that bound the tropical value of the Leibniz bound for $\mathcal N$ equals the notes'
  bound. Step~5 proves $2^{i+j+2-h}$ (one weaker), so the proved exponent exceeds the notes' by $O(h)$; the limit
  constant $\frac18\log2$ is the same. Lemma~\ref{lem:M} is used formally only at $q=5$.
\item \emph{Not formalised.} Proposition~\ref{prop:floor}, the remarks (including Remark~\ref{rem:173} on $q=173$)
  and Section~\ref{sec:limits}; none of them is used in the proofs of the theorems.
\end{itemize}

\subsection{How the certificates are checked}
A dual certificate is a record of rational data (for example \lean{Cert\slash{}DataQ14.lean},
\lean{TiQ\slash{}Cert\slash{}DataQ6L800.lean}, \lean{Cert\slash{}DataQ5L200.lean}, \lean{QZ\slash{}CertDataQ175.lean}): the log table as a list of integers in units of
$2^{-40}$, the heights $c_j$, the bounds $S_k$, $\lambda$ and $E_{\mathrm{lo}}$. A Boolean predicate (\lean{Cert.CData.Valid},
\lean{QZ.CDataZ.Valid}) collects every check of Section~\ref{sec:cert}, and the certificate files prove it by
\lean{decide +kernel}, i.e.\ by evaluation in the Lean kernel; soundness lemmas turn it into a \lean{DualCert}. The data were
produced by the scripts \lean{scripts\slash{}gen\_cert.py} and \lean{scripts\slash{}q2num\slash{}gen\_cert\_z.py}, which are not
trusted: a wrong value makes the build fail but cannot make a false statement provable. The block determinants of
Proposition~\ref{prop:blocks} are evaluated as rational functions of $q$ (\lean{Q\slash{}NonVan}, \lean{TiQ\slash{}NonVan\slash{}NonVanT})
and at $q/2$ (\lean{QZ\slash{}NonVan}) by exact computation.

\subsection{Building and replaying}
With elan installed:
\begin{quote}\ttfamily\small
git clone https://github.com/telaviv/chi2-irrationality\\
cd chi2-irrationality \&\& git checkout cef469e\\
lake exe cache get\\
LEAN\_NUM\_THREADS=4 lake build\\
lake env lean Axioms.lean
\end{quote}
The build has $9196$ jobs. Most of its time is spent in the kernel checks of the certificates: a few minutes each for
$W=1/14^2$, $W=1/6^2,\dots,1/12^2$ and $W=1/5^2$ ($461$\,s), about three hours for $W=4/175^2$ ($10{,}697$\,s wall time on a shared machine), and
several hours for the finer certificate at $W=1/36$, whose cells are split over four files. To avoid trusting the
elaborator's handling of the stored environment, each module can be replayed through the kernel with \lean{leanchecker}:
\begin{quote}\ttfamily\small
for m in \$(find Chi2Irrational -name "*.lean" | sed 's/\textbackslash.lean\$//; s\#/\#.\#g'); do\\
\hspace*{2em}lake env leanchecker "\$m"; done
\end{quote}
Each module is replayed against the compiled environment of its imports, so Mathlib itself is not re-checked. For the
formalisation of \cite{Chi2Paper} and the intermediate theorem (91 modules), all modules passed on 2026-11-09. For the
modules of Theorems~\ref{thm:B} and~\ref{thm:C} (\lean{TiQ\slash{}}, \lean{QZ\slash{}}), $40$ modules passed on 2026-11-12
($3189$\,s, log \lean{scripts\slash{}leanchecker\_qz\_tiq.log}); the seven modules that contain the \lean{decide +kernel}
proofs of the two certificates of Theorems~\ref{thm:B} ($q\ge6$) and~\ref{thm:C} (\lean{QZ\slash{}CertQ175.lean}, \lean{TiQ\slash{}Certs.lean} and the five files of
\lean{TiQ\slash{}Cert\slash{}} other than the data) were not replayed, because each would repeat hours of kernel evaluation;
they are checked by the kernel during \lean{lake build}. The modules added for Lemma~\ref{lem:M} and $q=5$
(\lean{TiQ\slash{}M\slash{}}, \lean{TiQ\slash{}Interface5.lean}, \lean{Cert\slash{}DataQ5L200.lean}, \lean{Cert\slash{}Q5L200.lean}) have not
been replayed.

\subsection{What has to be trusted}
Beyond the Lean kernel and the three axioms, the reader has to trust that the definitions displayed above express
Theorems~\ref{thm:A}--\ref{thm:C}, and the Mathlib definitions they use (real numbers, \lean{tsum}, \lean{Irrational},
\lean{Odd}). Nothing else, in particular no numerical data and no generating script, needs to be trusted. Apart from the
check by the Lean kernel, all verification of the formalisation so far was carried out by AI systems; no human expert has
yet reviewed the Lean code.

\section*{Acknowledgements and disclosure}
The research, the construction, the proofs, the certificate programs, the Lean formalisation, the literature searches and
the text of this paper were produced largely by AI models (Claude, Anthropic), working under the direction of the author.
Apart from the check by the Lean kernel, all verification so far, including the internal reviews of the argument by AI
referees, was carried out by AI systems. The written proofs of the new ingredients (the extension of the construction to
general $q$, Lemma~\ref{lem:tene}, the case of $\Ti_2$ including Lemma~\ref{lem:M}, and the case of $\chi_2(2/q)$) were
each checked by one AI referee,
at a single gate, with independent implementations of the machine. Several written proofs here follow the formal proofs
rather than those working notes (Lemma~\ref{lem:tene}, Lemma~\ref{lem:tenZ},
Proposition~\ref{prop:nonvanC}, and Theorem~\ref{thm:size} at $w<0$); as written proofs these were checked by the Lean kernel and by the whole-paper referee
only. Lemma~\ref{lem:M} (added after the whole-paper review) also follows the formal proof; as written it was checked by
the Lean kernel and by one AI referee of the changed sections, and its source proof by one AI gate referee. The fidelity of each Lean statement was checked by an independent AI referee, and the text of this paper was
reviewed as a whole by one further AI referee. No MathSciNet
search or expert query has been made. No human expert has yet reviewed the mathematics or the Lean code; we would be
grateful for corrections. The prime number theorem in the formalisation is ported from the PrimeNumberTheoremAnd project of
A.~Kontorovich, T.~Tao and collaborators \cite{PNTplus}, through the \lean{zeta5} repository of C\'esar Del Solar
\cite{zeta5}.

\bibliographystyle{amsplain}
\bibliography{references}

\end{document}



